% main.tex -- manuscript for the European Journal of Mechanics A/Solids, in Elsevier's
% CAS single-column class (cas-sc, not stored in the repository: build with
%   python paper/build.py --cls <folder with cas-sc.cls and cas-common.sty>).
\documentclass[a4paper,fleqn]{cas-sc}

\usepackage[numbers,sort&compress]{natbib}
\usepackage{amsthm}
% macros.tex -- notation of the manuscript.
%
% Every symbol of the paper goes through a semantic macro defined here, so that a
% change of notation is a one-line change.

% ---------------------------------------------------------------- packages
\usepackage{booktabs}
\usepackage{bm}
\usepackage{xcolor}
% hyperref and amsmath are loaded by the class (cas-sc)
% for paper/build.py, which checks that no word is set past the text block
\AtBeginDocument{\typeout{GRADINV-TEXTWIDTH=\the\textwidth}}

% ---------------------------------------------------------------- drafting aids
% \pending{...}: outline of what a section will say; printed in a grey box while
% \ifdraftnotes is true, invisible otherwise.
% \pendingcite{Key}: a sentence that depends on a source not read yet; prints a
% visible marker.
\newif\ifdraftnotes \draftnotestrue
\ifdefined\gradinvfinal\draftnotesfalse\fi  % set by paper/build.py --final
\newcommand{\pending}[1]{\ifdraftnotes\par\smallskip\noindent
  \fcolorbox{gray}{gray!8}{\parbox{\dimexpr\linewidth-2\fboxsep-2\fboxrule}{\small\sffamily
  \textbf{Pending.} #1}}\par\smallskip\fi}
\newcommand{\pendingcite}[1]{\textcolor{red}{[#1?]}}

% ---------------------------------------------------------------- fields
\newcommand{\disp}{\bm{u}}                 % displacement
\newcommand{\strain}{\bm{\varepsilon}}     % strain
\newcommand{\sgrad}{\bm{\eta}}             % strain gradient, form II: eta_ijk = eps_ij,k
\newcommand{\stress}{\bm{\sigma}}          % Cauchy-like stress
\newcommand{\dstress}{\bm{\tau}}           % double stress, tau_ijk = tau_jik

% ---------------------------------------------------------------- spaces and groups
\newcommand{\Vsg}{\mathcal{V}}             % strain-gradient space (dimension 18)
\newcommand{\End}{\operatorname{End}}
\newcommand{\Alg}[1]{\mathcal{A}_{#1}}     % commutant algebra of class #1
\newcommand{\Orth}{\mathrm{O}(3)}
\newcommand{\SOrth}{\mathrm{SO}(3)}
\newcommand{\GL}{\mathrm{GL}(3)}
\newcommand{\Hrm}[1]{H_{#1}}               % harmonic (irreducible O(3)) space of degree #1
\newcommand{\Hone}[1]{H_{1}^{#1}}          % copies of H_1: s (stretch) and r (rotation) of the GL(3) chain
\newcommand{\clsiso}{\mathrm{iso}}
\newcommand{\clscub}{m\bar{3}m}
\newcommand{\clshex}{6/mmm}

% ---------------------------------------------------------------- tensors
\newcommand{\tC}{\mathbb{C}}               % classical stiffness (fourth order)
\newcommand{\tD}{\mathbb{D}}               % gradient stiffness (sixth order), D_ijmkln
\newcommand{\tI}{\mathbb{I}}               % identity on V
\newcommand{\tK}{\mathbb{K}}               % generic element of an algebra
\newcommand{\tP}[1]{\mathbb{P}^{#1}}       % isotypic projector of irrep #1
\newcommand{\tU}[3]{\mathbb{U}^{#1}_{#2#3}} % basis unit: U^p_ab U^q_cd = delta_pq delta_bc U^p_ad
\newcommand{\tOct}{\mathbb{O}}            % cubic structure tensor sum_s e_s^{(x)4} (delta_ijkl of Lazar et al. 2022); components O_ijkl written blackboard (italic O is the group)
\newcommand{\tHsix}{\mathbb{H}_{6}}        % hexagonal structure tensor Re(a + i n x a)^{(x)6}
\newcommand{\tB}{\mathbb{B}}               % concentration tensor of the self-consistent scheme
\newcommand{\tS}{\mathbb{S}}               % Eshelby tensors (application, appendix B)
\newcommand{\tSav}{\bar{\mathbb{S}}}       % sphere average of the sixth-order Eshelby tensor (Ma's <N^>, Solyaev's <S^(6)>_v)
\newcommand{\gr}{\mathrm{g}}               % subscript of the grain tensors in section 7
\newcommand{\tdot}{\mathbin{\vdots}}       % triple contraction
\newcommand{\ot}{\otimes}
\newcommand{\kron}{\bm{\delta}}            % second-order identity

\newcommand{\emb}[1]{J_{#1}}               % isometric equivariant embedding of a copy: J^p_a (E is an irrep label)

\newcommand{\harm}{\mathcal{H}}            % harmonic (totally symmetric traceless) part of a third-order tensor
\newcommand{\zhex}{\bm{z}}                 % complex vector a + i n x a of the hexagonal structure tensors

% ---------------------------------------------------------------- blocks
\newcommand{\blk}[2]{#1^{#2}}              % block of tensor #1 in irrep #2, e.g. \blk{M}{T_1}
\newcommand{\hexA}{\mathsf{A}}             % hexagonal A_2 block (3 x 3), copies (3, s, r)
\newcommand{\hexF}{\mathsf{F}}             % hexagonal E_1 block (4 x 4), copies (3, 2, s, r)
\newcommand{\hexG}{\mathsf{G}}             % hexagonal E_2 block (2 x 2), copies (3, 2)
\newcommand{\irr}[1]{#1}                   % Mulliken label of an irrep (A_2, E, T_1, ...)
\newcommand{\Dhex}{\Delta}                 % hexagonal extra modulus, m_{B_1} - m_{B_2}
\newcommand{\dimirr}[1]{d_{#1}}            % dimension of irrep #1
\newcommand{\mult}[1]{n_{#1}}              % multiplicity of irrep #1 in V

% ---------------------------------------------------------------- averages and lengths
\newcommand{\Vg}[1]{#1^{\mathrm{V}}}       % Voigt
\newcommand{\Rs}[1]{#1^{\mathrm{R}}}       % Reuss
\newcommand{\Hl}[1]{#1^{\mathrm{H}}}       % Hill
\newcommand{\avg}[1]{\langle #1\rangle}    % orientation average
\newcommand{\lblk}[1]{\ell_{#1}}           % block length
\newcommand{\lLP}{\ell_{\mathrm{LP}}}      % Lazar-Po length
\newcommand{\linf}{\ell_{\infty}}          % large-grain limit of the self-consistent length
\newcommand{\scres}{\rho}                  % residual of the self-consistent condition (was Delta)

% ---------------------------------------------------------------- misc
\newcommand{\dd}{\mathrm{d}}
\newcommand{\iu}{\mathrm{i}}
\newcommand{\tr}{\operatorname{tr}}
\newcommand{\sym}{\operatorname{sym}}
\newcommand{\angstrom}{\text{\normalfont\AA}}  % unit in math mode

% Generated by analysis/make_numbers.py -- do not edit by hand.
% Each macro: value, then its provenance.
\newcommand{\NumAlmE}{0.03560} % cubic.from_lazar2022
\newcommand{\NumAlmA}{0.93596} % cubic.from_lazar2022
\newcommand{\NumCuMEAMmE}{0.00005} % cubic.from_lazar2022; rounding of the inputs +-2e-06 (rough)
\newcommand{\NumCuMEAMmA}{0.22441} % cubic.from_lazar2022
\newcommand{\NumCuEAMmE}{0.08546} % cubic.from_lazar2022
\newcommand{\NumCuEAMmA}{0.72458} % cubic.from_lazar2022
\newcommand{\NumCuMEAMLambdaSeven}{0.6286} % cubic.from_lazar2022 (eigenvalues of the T1 block)
\newcommand{\NumFeNthreethreeT}{-0.238} % cubic.from_lazar2022
\newcommand{\NumFeNthreetwoT}{-0.470} % cubic.from_lazar2022
\newcommand{\NumFeDetNT}{-0.318} % cubic.from_lazar2022
\newcommand{\NumWNthreethreeT}{-0.022} % cubic.from_lazar2022
\newcommand{\NumWNthreetwoT}{-1.347} % cubic.from_lazar2022
\newcommand{\NumWDetNT}{-1.836} % cubic.from_lazar2022
\newcommand{\NumFeNeighborhoodRadius}{1.76} % lengths.neighborhood_radius (A)
\newcommand{\NumFeNeighborhoodRatio}{0.71} % lengths.neighborhood_radius / b
\newcommand{\NumWNeighborhoodRadius}{1.52} % lengths.neighborhood_radius (A)
\newcommand{\NumWNeighborhoodRatio}{0.55} % lengths.neighborhood_radius / b
\newcommand{\NumAlVoigtmthree}{0.73882} % averages.voigt
\newcommand{\NumAlVoigtmtwo}{0.08986} % averages.voigt; rounding of the inputs +-2e-05
\newcommand{\NumAlVoigtMss}{1.38931} % averages.voigt
\newcommand{\NumAlVoigtMsr}{0.40536} % averages.voigt
\newcommand{\NumAlVoigtMrr}{0.94719} % averages.voigt
\newcommand{\NumAlReussmthree}{0.34346} % averages.reuss
\newcommand{\NumAlReussmtwo}{0.04497} % averages.reuss; rounding of the inputs +-2e-05
\newcommand{\NumAlReussMss}{1.38922} % averages.reuss
\newcommand{\NumAlReussMsr}{0.41141} % averages.reuss
\newcommand{\NumAlReussMrr}{0.5034} % averages.reuss; rounding of the inputs +-5e-05
\newcommand{\NumAlReussRatio}{1.263} % averages.reuss_mE_ratio
\newcommand{\NumCuMEAMVoigtmthree}{0.46233} % averages.voigt
\newcommand{\NumCuMEAMVoigtmtwo}{0.02986} % averages.voigt; rounding of the inputs +-1e-05
\newcommand{\NumCuMEAMVoigtMss}{1.46428} % averages.voigt
\newcommand{\NumCuMEAMVoigtMsr}{0.36030} % averages.voigt
\newcommand{\NumCuMEAMVoigtMrr}{0.77210} % averages.voigt
\newcommand{\NumCuMEAMReussmthree}{0.15496} % averages.reuss
\newcommand{\NumCuMEAMReussmtwo}{0.000125} % averages.reuss; rounding of the inputs +-5e-06 (rough)
\newcommand{\NumCuMEAMReussMss}{0.9296} % averages.reuss; rounding of the inputs +-6e-05
\newcommand{\NumCuMEAMReussMsr}{-0.0149} % averages.reuss; rounding of the inputs +-5e-05
\newcommand{\NumCuMEAMReussMrr}{0.5088} % averages.reuss; rounding of the inputs +-5e-05
\newcommand{\NumCuMEAMReussRatio}{2.4960} % averages.reuss_mE_ratio
\newcommand{\NumCuEAMVoigtmthree}{0.63528} % averages.voigt
\newcommand{\NumCuEAMVoigtmtwo}{0.14265} % averages.voigt; rounding of the inputs +-2e-05
\newcommand{\NumCuEAMVoigtMss}{2.1640} % averages.voigt; rounding of the inputs +-5e-05
\newcommand{\NumCuEAMVoigtMsr}{-0.05317} % averages.voigt
\newcommand{\NumCuEAMVoigtMrr}{0.23777} % averages.voigt
\newcommand{\NumCuEAMReussmthree}{0.54107} % averages.reuss
\newcommand{\NumCuEAMReussmtwo}{0.12459} % averages.reuss; rounding of the inputs +-2e-05
\newcommand{\NumCuEAMReussMss}{1.8637} % averages.reuss; rounding of the inputs +-8e-05
\newcommand{\NumCuEAMReussMsr}{-0.03191} % averages.reuss
\newcommand{\NumCuEAMReussMrr}{0.23626} % averages.reuss
\newcommand{\NumCuEAMReussRatio}{1.458} % averages.reuss_mE_ratio
\newcommand{\NumFeBonePrinted}{7.29} % selfconsistent.invariants
\newcommand{\NumFeBoneMa}{11.56} % selfconsistent.invariants
\newcommand{\NumFeBoneExact}{6.61} % selfconsistent.invariants
\newcommand{\NumFeSixthExact}{0.65} % selfconsistent.invariants
\newcommand{\NumFeCoherentEllTen}{1.001} % selfconsistent.solve
\newcommand{\NumFeCoherentNuTen}{0.3483} % selfconsistent.solve
\newcommand{\NumFeCoherentResidualTen}{10^{-4}} % selfconsistent.solve
\newcommand{\NumFeCoherentEllFifty}{1.001} % selfconsistent.solve
\newcommand{\NumFeCoherentNuFifty}{0.3485} % selfconsistent.solve
\newcommand{\NumFeCoherentResidualFifty}{4\times10^{-5}} % selfconsistent.solve
\newcommand{\NumSixthMin}{-0.649} % test_selfconsistent
\newcommand{\NumSixthMax}{0.410} % test_selfconsistent
\newcommand{\NumSixthMaxPercent}{22} % test_selfconsistent
\newcommand{\NumAlLimEllSeven}{1.388} % selfconsistent.large_grain_limit
\newcommand{\NumAlLimEllEight}{1.378} % selfconsistent.large_grain_limit
\newcommand{\NumAlLimResidualSeven}{0.26} % selfconsistent.large_grain_limit
\newcommand{\NumAlKronerNu}{0.3356} % selfconsistent.kroner_nu
\newcommand{\NumAlLengthCrystal}{1.177} % Lazar-Po eq. 107
\newcommand{\NumAlLengthVoigt}{1.170} % Lazar-Po eq. 107
\newcommand{\NumAlLengthReuss}{0.883} % Lazar-Po eq. 107
\newcommand{\NumCuMEAMLimEllSeven}{1.082} % selfconsistent.large_grain_limit
\newcommand{\NumCuMEAMLimEllEight}{1.074} % selfconsistent.large_grain_limit
\newcommand{\NumCuMEAMLimResidualSeven}{0.40} % selfconsistent.large_grain_limit
\newcommand{\NumCuMEAMKronerNu}{0.3360} % selfconsistent.kroner_nu
\newcommand{\NumCuMEAMLengthCrystal}{0.785} % Lazar-Po eq. 107
\newcommand{\NumCuMEAMLengthVoigt}{0.753} % Lazar-Po eq. 107
\newcommand{\NumCuMEAMLengthReuss}{0.565} % Lazar-Po eq. 107
\newcommand{\NumCuEAMLimEllSeven}{1.492} % selfconsistent.large_grain_limit
\newcommand{\NumCuEAMLimEllEight}{1.475} % selfconsistent.large_grain_limit
\newcommand{\NumCuEAMLimResidualSeven}{0.60} % selfconsistent.large_grain_limit
\newcommand{\NumCuEAMKronerNu}{0.3392} % selfconsistent.kroner_nu
\newcommand{\NumCuEAMLengthCrystal}{0.934} % Lazar-Po eq. 107
\newcommand{\NumCuEAMLengthVoigt}{0.907} % Lazar-Po eq. 107
\newcommand{\NumCuEAMLengthReuss}{0.841} % Lazar-Po eq. 107
\newcommand{\NumTableSixMaxDiff}{7\times10^{-6}} % cubic_results; Lazar et al. 2022 Table 6
\newcommand{\NumAlCond}{47.9} % cubic_results
\newcommand{\NumAlRVmthree}{0.4649} % cubic_results (Reuss/Voigt)
\newcommand{\NumAlRVmtwo}{0.5005} % cubic_results (Reuss/Voigt)
\newcommand{\NumAlRVHone}{0.4603} % cubic_results (Reuss/Voigt, H1)
\newcommand{\NumAlVoigtLengthSpread}{2.91} % lengths.block_lengths
\newcommand{\NumAlHillLengthLP}{1.037} % lengths.lazar_po_length
\newcommand{\NumCuMEAMCond}{3.4\times10^{4}} % cubic_results; rounding of the inputs +-1e+03 (rough)
\newcommand{\NumCuMEAMRVmthree}{0.3352} % cubic_results (Reuss/Voigt)
\newcommand{\NumCuMEAMRVmtwo}{0.0042} % cubic_results (Reuss/Voigt); rounding of the inputs +-0.0002 (rough)
\newcommand{\NumCuMEAMRVHone}{0.4724} % cubic_results (Reuss/Voigt, H1)
\newcommand{\NumCuMEAMVoigtLengthSpread}{5.15} % lengths.block_lengths
\newcommand{\NumCuMEAMHillLengthLP}{0.671} % lengths.lazar_po_length
\newcommand{\NumCuEAMCond}{26.2} % cubic_results
\newcommand{\NumCuEAMRVmthree}{0.8517} % cubic_results (Reuss/Voigt)
\newcommand{\NumCuEAMRVmtwo}{0.8734} % cubic_results (Reuss/Voigt)
\newcommand{\NumCuEAMRVHone}{0.8585} % cubic_results (Reuss/Voigt, H1)
\newcommand{\NumCuEAMVoigtLengthSpread}{3.69} % lengths.block_lengths
\newcommand{\NumCuEAMHillLengthLP}{0.882} % lengths.lazar_po_length
\newcommand{\NumAlLimResidualEight}{0.30} % application.row (selfconsistent.large_grain_limit)
\newcommand{\NumAlVoigtKronerSpread}{2.91} % application.row (lengths.block_lengths, Kroner C)
\newcommand{\NumAlLimOverLengthPercent}{19} % application.row (l_inf, 7 terms, over l_LP of Voigt with Kroner C)
\newcommand{\NumAlLimShiftPercent}{0.7} % application.row (l_inf, 7 to 8 terms)
\newcommand{\NumCuMEAMLimResidualEight}{0.43} % application.row (selfconsistent.large_grain_limit)
\newcommand{\NumCuMEAMVoigtKronerSpread}{5.10} % application.row (lengths.block_lengths, Kroner C)
\newcommand{\NumCuMEAMLimOverLengthPercent}{44} % application.row (l_inf, 7 terms, over l_LP of Voigt with Kroner C)
\newcommand{\NumCuMEAMLimShiftPercent}{0.7} % application.row (l_inf, 7 to 8 terms)
\newcommand{\NumCuEAMLimResidualEight}{0.66} % application.row (selfconsistent.large_grain_limit)
\newcommand{\NumCuEAMVoigtKronerSpread}{3.87} % application.row (lengths.block_lengths, Kroner C)
\newcommand{\NumCuEAMLimOverLengthPercent}{65} % application.row (l_inf, 7 terms, over l_LP of Voigt with Kroner C)
\newcommand{\NumCuEAMLimShiftPercent}{1.1} % application.row (l_inf, 7 to 8 terms)
\newcommand{\NumSCCases}{72} % selfconsistent_fcc.json (analysis/selfconsistent_fcc.py)
\newcommand{\NumSCNearLimitPercent}{2} % selfconsistent_fcc.json (analysis/selfconsistent_fcc.py)
\newcommand{\NumSCSixthSolutionMaxPercent}{4} % selfconsistent_fcc.json (analysis/selfconsistent_fcc.py)
\newcommand{\NumSCResidualTwoNmMin}{0.018} % selfconsistent_fcc.json (analysis/selfconsistent_fcc.py)
\newcommand{\NumSCResidualTwoNmMax}{0.044} % selfconsistent_fcc.json (analysis/selfconsistent_fcc.py)
\newcommand{\NumSCResidualHundredNmMin}{3.6\times10^{-4}} % selfconsistent_fcc.json (analysis/selfconsistent_fcc.py)
\newcommand{\NumSCResidualHundredNmMax}{9.0\times10^{-4}} % selfconsistent_fcc.json (analysis/selfconsistent_fcc.py)
\newcommand{\NumCuMEAMSensitivityEll}{10^{-7}} % selfconsistent_fcc.json (analysis/selfconsistent_fcc.py)
\newcommand{\NumAlInducedCrossoverNm}{5.6} % induced_moduli.summary (Malyi 2021, eqs. 10.21-10.22)
\newcommand{\NumCuMEAMInducedCrossoverNm}{0.7} % induced_moduli.summary (Malyi 2021, eqs. 10.21-10.22)
\newcommand{\NumCuEAMInducedCrossoverNm}{0.8} % induced_moduli.summary (Malyi 2021, eqs. 10.21-10.22)
\newcommand{\NumAlInducedRatioHundredNm}{3\times10^{2}} % induced_moduli.summary (Malyi 2021, eqs. 10.21-10.22)
\newcommand{\NumCuMEAMInducedRatioHundredNm}{2\times10^{4}} % induced_moduli.summary (Malyi 2021, eqs. 10.21-10.22)
\newcommand{\NumCuEAMInducedRatioHundredNm}{2\times10^{4}} % induced_moduli.summary (Malyi 2021, eqs. 10.21-10.22)
\newcommand{\NumFeZener}{1.23} % fe_case.json (analysis/fe_case.py)
\newcommand{\NumFeMuVoigt}{26.18} % fe_case.json (analysis/fe_case.py)
\newcommand{\NumFeMuReuss}{25.91} % fe_case.json (analysis/fe_case.py)
\newcommand{\NumFeKronerNu}{0.3485} % fe_case.json (analysis/fe_case.py)
\newcommand{\NumFeLimEll}{1.0012} % fe_case.json (analysis/fe_case.py)
\newcommand{\NumFeExactEllMin}{1.0007} % fe_case.json (analysis/fe_case.py)
\newcommand{\NumFeExactEllMax}{1.0011} % fe_case.json (analysis/fe_case.py)
\newcommand{\NumFeExactNuMin}{0.3483} % fe_case.json (analysis/fe_case.py)
\newcommand{\NumFeExactNuMax}{0.3485} % fe_case.json (analysis/fe_case.py)
\newcommand{\NumFeExactResidualMax}{2\times10^{-4}} % fe_case.json (analysis/fe_case.py)
\newcommand{\NumFeAllEllMin}{1.0001} % fe_case.json (analysis/fe_case.py)
\newcommand{\NumFeAllEllMax}{1.0013} % fe_case.json (analysis/fe_case.py)
\newcommand{\NumFeAllResidualMax}{5\times10^{-4}} % fe_case.json (analysis/fe_case.py)
\newcommand{\NumFeAnisoZener}{2.46} % fe_case.json (analysis/fe_case.py)
\newcommand{\NumFeAnisoEllMin}{1.013} % fe_case.json (analysis/fe_case.py)
\newcommand{\NumFeAnisoEllMax}{1.022} % fe_case.json (analysis/fe_case.py)
\newcommand{\NumFeMixedCount}{12} % fe_case.json (analysis/fe_case.py)
\newcommand{\NumFeMixedRmsMin}{0.16} % fe_case.json (analysis/fe_case.py)
\newcommand{\NumFeMixedRmsMax}{0.57} % fe_case.json (analysis/fe_case.py)
\newcommand{\NumHexIsoReussAtSigma}{0.9130} % hexagonal_granular
\newcommand{\NumHexTIReussAtSigma}{0.9083} % hexagonal_granular
\newcommand{\NumHexGranDeltaMax}{0.21} % hexagonal_granular
\newcommand{\NumHexGranReussChangePercent}{0.4} % hexagonal_granular
\newcommand{\NumCuMEAMmERoundingPercent}{4} % claim_audit.uncertainties


\newtheorem{proposition}{Proposition}[section]
\newtheorem{lemma}[proposition]{Lemma}
\theoremstyle{remark}
\newtheorem{remark}[proposition]{Remark}

% E-mail address as text, without the class icon file (thumbnails/cas-email.jpeg):
% Editorial Manager cannot process subfolders.
\ExplSyntaxOn
\keys_set:nn { stm / mktitle } { nologo }
\ExplSyntaxOff

\begin{document}
\let\WriteBookmarks\relax
\def\floatpagepagefraction{1}
\def\textpagefraction{.001}

\shorttitle{Algebra and averages of cubic and hexagonal strain-gradient tensors}
\shortauthors{D. Gonz\'alez-Monsalve}

\title[mode=title]{Algebra, closed-form compliance and orientation averages of cubic and
hexagonal strain-gradient elasticity tensors}

\author[1]{Daniel Gonz\'alez-Monsalve}[orcid=0000-0002-3892-4722]
\cormark[1]
\ead{danielgonzalezm@udec.cl}

\affiliation[1]{organization={Department of Mechanical Engineering, Faculty of Engineering,
  Universidad de Concepci\'on},
  addressline={Edmundo Larenas 219},
  postcode={4070409},
  city={Concepci\'on},
  country={Chile}}

\cortext[1]{Corresponding author}

\begin{abstract}
% at most 250 words (EJM A/Solids), no references or undefined abbreviations;
% counted by tests/test_manuscript.py
Strain-gradient elasticity tensors of crystals are computed from interatomic
potentials, but closed-form algebras for their products, inverses and averages have been
published only for isotropy and transverse isotropy. We construct coordinate-free bases,
with closed multiplication tables, of the algebras of cubic ($m\bar3m$) and hexagonal
($6/mmm$) sixth-order tensors on the strain-gradient space of Mindlin's form~II; their
symmetric elements are the gradient tensors, with 11 and 22 constants, and every operation
reduces to blocks of size at most four. The hexagonal modulus absent from
transverse isotropy is proportional to the component of the tensor along the sixth-order
hexagonal structure tensor; in a granular packing it measures the sixfold harmonic of the
contact distribution. We give in closed form the gradient compliance, with block-wise
positivity conditions and changes of basis to published parametrisations, and the
Voigt, Reuss and Hill estimates for untextured polycrystals, with Reuss below Voigt except
for isotropic grains. These estimates omit the moduli induced by grain heterogeneity,
which dominate beyond a few nanometres. The formulas are applied to atomistic data of
cubic metals and illustrated with a self-consistent scheme, correcting a published sphere
average of the sixth-order Eshelby tensor.
\end{abstract}

% 1 to 7 keywords (EJM A/Solids); the highlights go in a separate file at submission
\begin{keywords}
strain-gradient elasticity \sep sixth-order tensors \sep cubic symmetry \sep
hexagonal symmetry \sep tensor algebra \sep orientation averages \sep polycrystals
\end{keywords}

\maketitle

\section{Introduction}\label{sec:intro}

Strain-gradient elasticity adds to the classical energy a quadratic form in the strain
gradient, governed in Mindlin's form~II by a sixth-order tensor
$\tD$~\cite{Mindlin1964,MindlinEshel1968}. For crystals it is no longer unknown: Admal et al.~\cite{Admal2017} derived its atomistic
representation, and Lazar et al.~\cite{Lazar2022} and Po et al.~\cite{Po2019} published
complete cubic tensors of fcc and bcc metals from interatomic potentials, which now enter non-singular theories of dislocations~\cite{Po2018,Lazar2024};
granular micromechanics gives anisotropic gradient tensors in terms of the
fabric~\cite{Pirmoradi2024}. Using such tensors for heterogeneous bodies
requires the operations that classical micromechanics performs on
fourth-order tensors: products, inverses, orientation averages and concentration tensors.
Inverses of sixth-order tensors appear in the inclusion and inhomogeneity problems of
Monchiet and Bonnet~\cite{Monchiet2011a,Monchiet2013,Monchiet2015} and of Ma et al.~\cite{Ma2018},
the compliance in estimates of Reuss type~\cite{Nazarenko2021a}, and all four operations
in the self-consistent scheme of Solyaev~\cite{Solyaev2023}.

For fourth-order tensors, Walpole~\cite{Walpole1984} gave bases with multiplication tables
for the 32 crystal classes, which reduce these operations to products of small matrices;
bases of the same kind also give the orientation averages of third-order elastic
constants~\cite{Kube2026}, and harmonic bases of tensors of any order, with tensorial
texture coefficients, follow from the representations of the rotation
group~\cite{Krause2025}. For gradient tensors, the corresponding algebra is known for
isotropy~\cite{Monchiet2011a,Auffray2015}, also without the minor
symmetries~\cite{Monchiet2011b}, and for transverse isotropy~\cite{Monchiet2013}, where
it was built to solve inhomogeneity and crack problems~\cite{Monchiet2015}. The anisotropic tensors have been classified. Olive and Auffray~\cite{Olive2013}
obtained the 17 symmetry classes of the sixth-order tensor, and Auffray et
al.~\cite{Auffray2013} wrote the matrix of each class in an orthonormal basis, and later
for the 48 classes of the whole law, fifth-order coupling included~\cite{Auffray2019}. Le Quang
et al.~\cite{LeQuang2021} classified the tensors by their symmetry planes, Auffray et
al.~\cite{Auffray2021} made the harmonic structure explicit in two dimensions, and
Poorsolhjouy~\cite[table~B.2]{Poorsolhjouy2026} counted, class by class, the tensors
without major symmetry. Demonteaux et al.~\cite{Demonteaux2025} report an orthogonal basis for the cubic
tensor, obtained from the character and projection formulas of the cubic group; its
formulas are not yet published. For the cubic class, Lazar et al.~\cite{Lazar2022} wrote
the tensor with eleven constants, reduced its matrix to diagonal blocks of sizes five and
three, and obtained its eigenvalues by factoring the characteristic polynomial of the
block of size five into a quadratic and a cubic, conditions of positive definiteness and Voigt
average. Averages of the same Voigt type have been computed for fcc metals in Mindlin's
second-gradient theory~\cite{Bagherpour2024a,Bagherpour2024b}. Nazarenko et
al.~\cite{Nazarenko2021a} gave the inverse of the coupled constitutive law in tensorial
form, explicitly for hemitropic materials.

These works establish how many constants each class has and where they sit in a matrix:
they are parametrisations. An algebra is a different object, a basis in which the
product of two tensors of the class is a product of small blocks, so that the inverse,
the orientation average and the concentration tensor follow block by block. To our
knowledge no such basis, and no closed-form compliance or Reuss estimate, has been
published for the cubic or the hexagonal class. For the cubic class an orthogonal basis
has been announced~\cite{Demonteaux2025}; what we add there are its explicit formulas,
the multiplication table and their uses. The hexagonal class has a further interest.
In two dimensions, the gradient tensor of a hexagonal lattice is anisotropic whereas its
classical tensor is isotropic, which makes wave propagation anisotropic at high
frequency~\cite{Rosi2016}. In three dimensions the hexagonal classical tensor is
transversely isotropic, and the single gradient modulus that breaks transverse isotropy
has been identified only as the coefficient of a fixed matrix~\cite{Auffray2013}, which
we write $\eta_{\mathrm A}$ (their $\eta$) to keep $\sgrad$ for the strain gradient.

This paper constructs these algebras and uses them. First, we give coordinate-free bases,
with closed multiplication tables, of the algebras of the classes $\clscub$ and $\clshex$
on the strain-gradient space, of dimensions 15 and 32. Their symmetric elements are the
gradient tensors, with 11 and 22 constants, and the other elements are the
non-reciprocal tensors of the same classes. Repeated irreducible components are labelled
by the chain $\GL\supset\SOrth\supset G$, with the conversion to the labels of
Lazar~\cite{Lazar2016}. The hexagonal modulus absent from transverse isotropy is the
splitting of the part of azimuthal order $\pm3$ of the totally symmetric harmonic
component into two irreducible parts, equivalently, up to a factor, the component of
$\tD$ along the hexagonal structure tensor $\mathrm{Re}\,(\bm a+\iu\,\bm n\times\bm a)^{\ot6}$, and it
equals $4\eta_{\mathrm A}$. The hexagonal algebra is the transversely isotropic algebra of
Monchiet and Bonnet~\cite{Monchiet2013} with one of its idempotents split in two; what is
new there is the orthonormal basis, the identification of the extra modulus and the closed
forms. Second, the gradient compliance $\tD^{-1}$ follows in closed block form,
with conditions of invertibility and positive definiteness block by block and the change of basis to and from the constants of Lazar et
al.~\cite{Lazar2022} and the matrices of Auffray et al.~\cite{Auffray2013}. Third, the
Voigt, Reuss and Hill estimates of untextured polycrystals are weighted means of the
block moduli. The Reuss estimate lies below the Voigt one, with equality only for an
isotropic crystal; in the cubic case their vector-like blocks differ by a matrix of rank
one, and the hexagonal extra modulus drops out of the Voigt estimate. The cubic formulas
are applied to the published atomistic data of fcc and bcc metals. Finally, as an illustration, the self-consistent
scheme of Solyaev~\cite{Solyaev2023} with the complete cubic tensor reduces to the
inversion of four blocks and, for grains of physical size, sees the grain only through
its Voigt estimate. On the way we correct the average over the sphere of the sixth-order
Eshelby tensor of Ma et al.~\cite{Ma2018}, on which that scheme relies.

The averages are algebraic estimates, not bounds, and the paper is not a homogenisation
theory: the heterogeneity of the classical stiffness alone induces gradient moduli that
grow with the square of the grain size~\cite{Malyi2021}, and no average of $\tD$ contains
them. By Malyi's estimate they equal the Voigt estimate of $\tD$ at a grain diameter of
about $\NumAlInducedCrossoverNm$~nm for Al and $\NumCuMEAMInducedCrossoverNm$~nm for Cu
and dominate above it (section~\ref{sec:averages-cubic}). The averages are therefore the
part of the effective tensor carried by $\tD$: an isotropic substitute for the crystal
tensor, as in~\cite[section~7]{Lazar2022}, the gradient tensor of nanocrystals whose grains
are a few lattice parameters across, and an ingredient of a scheme that also accounts for
the heterogeneity of $\tC$. No hexagonal gradient tensor of a crystal has been published, so the hexagonal
results are illustrated with the granular micromechanics of Pirmoradi
et al.~\cite{Pirmoradi2024}, in which the extra modulus is the sixfold harmonic of the
contact distribution.

Section~\ref{sec:setting} recalls form~II and builds the isotropic algebra as a template,
sections~\ref{sec:cubic}--\ref{sec:application} treat the cubic and hexagonal algebras,
the compliance with the data and the verification, the averages and the illustration,
and the appendices convert our labels to those of Lazar and derive the averaged Eshelby
tensor.

\section{Strain-gradient elasticity, symmetry classes and the isotropic algebra}\label{sec:setting}

\subsection{Form II and the strain-gradient space}\label{sec:form-II}

In Mindlin's form~II~\cite{Mindlin1964,MindlinEshel1968}, the energy density is a
quadratic function of the strain $\varepsilon_{ij}=u_{(i,j)}$ and of the strain gradient,
written with the derivative index last as in~\cite{Lazar2016,Lazar2022},
$\eta_{ijk}=\varepsilon_{ij,k}$:
\begin{equation}
W=\tfrac12\,\strain:\tC:\strain+\tfrac12\,\sgrad\tdot\tD\tdot\sgrad,\qquad
\stress=\tC:\strain,\qquad \dstress=\tD\tdot\sgrad,\qquad
\tau_{ijm}=D_{ijmkln}\,\eta_{kln},\label{eq:energy}
\end{equation}
with $D_{ijmkln}=D_{jimkln}=D_{ijmlkn}=D_{klnijm}$. The fifth-order tensor coupling
$\strain$ and $\sgrad$ vanishes for centrosymmetric materials~\cite{Lazar2022} and is
absent throughout.

The strain gradients form the space $\Vsg$ of third-order tensors symmetric in their first
two indices (dimension 18, inner product $\eta_{ijk}\eta'_{ijk}$), and $\tD$ is a linear
map of $\Vsg$. Composition is written by juxtaposition,
$(\mathbb A\mathbb B)_{ijmkln}=A_{ijmpqr}B_{pqrkln}$; the identity is
$I_{ijmkln}=\tfrac12(\delta_{ik}\delta_{jl}+\delta_{il}\delta_{jk})\delta_{mn}$, the
transpose $A^{\mathsf T}_{ijmkln}=A_{klnijm}$ and the inner product
$\langle\mathbb A,\mathbb B\rangle=A_{ijmkln}B_{ijmkln}$; major symmetry is
$\tD=\tD^{\mathsf T}$. Form~I uses $\kappa_{ijk}=u_{i,jk}$, of which $\sgrad$ is a linear
combination~\cite{MindlinEshel1968}; with the isotropic, invertible
but not isometric map $\mathbb T:\kappa\mapsto\sgrad$,
$\eta_{ijk}=\tfrac12(\kappa_{ijk}+\kappa_{jik})$, the form-I tensor is
$\mathbb T^{\mathsf T}\tD\,\mathbb T$. The orthogonal bases below belong to form~II;
form~I enters only through the sources, in this index order for the granular tensors of
section~\ref{sec:averages-hexagonal} and as $\tilde\kappa_{ijk}=\kappa_{kij}=u_{k,ij}$ in
section~\ref{sec:application}.

\subsection{Symmetry classes and their algebras}\label{sec:algebras}

A rotation $Q$ acts on $\Vsg$ by the orthogonal map
$(Q\star\sgrad)_{ijk}=Q_{ip}Q_{jq}Q_{kr}\eta_{pqr}$. The inversion acts as $-1$ and
leaves every linear map of $\Vsg$ unchanged, so the symmetry of $\tD$ is fixed by a
subgroup $G$ of $\SOrth$, and the algebra of the class is the commutant
\begin{equation}
\Alg{G}=\{\tK\in\End(\Vsg):\ (Q\star)\,\tK=\tK\,(Q\star)\ \text{for all } Q\in G\}.
\label{eq:commutant}
\end{equation}
It is closed under composition, transposition and inversion. The tensors with major
symmetry are its symmetric elements, a subspace but not a subalgebra; the others are the
non-reciprocal (`odd') tensors of the class~\cite{Poorsolhjouy2026}. We write
$\Alg{\clsiso}$, $\Alg{\clscub}$ and $\Alg{\clshex}$ for $G=\SOrth$, the octahedral group
$O$ and $D_6$: crystals of point groups $\clscub$ and $\clshex$, which contain these
rotations and the inversion, have gradient tensors of these classes and no fifth-order
coupling. Lattices with a basis, such as diamond and hcp, are outside the scope, though
not because of the class of $\tD$. Their atomistic tensors~\cite{Admal2017} are
referred to a chosen reference point of the lattice, and the fifth-order coupling depends
on that choice. At an atom site, of point group $\bar43m$ or $\bar6m2$, the improper
operations act on $\tD$ as rotations, so $\tD$ is of class $\clscub$ or $\clshex$, but
without the inversion the fifth-order coupling does not vanish. At an inversion centre
it vanishes~\cite[theorem~1]{Admal2017}, but the point group is only $\bar3m$.

Let the irreducible representation $p$ of $G$, of dimension $\dimirr{p}$, occur
$\mult{p}$ times in $\Vsg$. Fix a model space $\mathcal W_p$ and, for each copy
$a=1,\dots,\mult{p}$, an isometric $G$-equivariant embedding $\emb{a}^{p}:\mathcal
W_p\to\Vsg$, with orthogonal images. Every irreducible representation of $\SOrth$, $O$
and $D_6$ is of real type, so by Schur's lemma the units
$\tU{p}{a}{b}=\emb{a}^{p}(\emb{b}^{p})^{\mathsf T}$ form a basis of $\Alg{G}$, with
\begin{equation}
\tU{p}{a}{b}\,\tU{q}{c}{d}=\delta_{pq}\delta_{bc}\,\tU{p}{a}{d},\qquad
(\tU{p}{a}{b})^{\mathsf T}=\tU{p}{b}{a},\qquad
\langle\tU{p}{a}{b},\tU{q}{c}{d}\rangle=\dimirr{p}\,\delta_{pq}\delta_{ac}\delta_{bd}.
\label{eq:table}
\end{equation}
An element is $\tK=\sum_p\sum_{a,b}K^p_{ab}\,\tU{p}{a}{b}$, with one block
$K^p\in\mathbb R^{\mult{p}\times\mult{p}}$ per irreducible representation:
$\Alg{G}\cong\bigoplus_p\mathbb R^{\mult{p}\times\mult{p}}$ has dimension
$\sum_p\mult{p}^2$, of which $\sum_p\tfrac12\mult{p}(\mult{p}+1)$ are symmetric
(table~\ref{tab:algebras}). The table~\eqref{eq:table} reduces every operation to the
blocks:
\begin{equation}
(\tK\mathbb L)^p=K^pL^p,\qquad
\tK^{-1}=\sum_p\sum_{a,b}\big[(K^p)^{-1}\big]_{ab}\,\tU{p}{a}{b},\qquad
K^p_{ab}=\langle\tK,\tU{p}{a}{b}\rangle/\dimirr{p}.\label{eq:block-rules}
\end{equation}
Hence $\tK$ is invertible if and only if every block is, the eigenvalues of a symmetric
$\tK$ are those of its blocks, each $\dimirr{p}$-fold, and $\tK$ is positive definite if
and only if every block is. For any linear map of $\Vsg$, the last formula
of~\eqref{eq:block-rules} gives its orthogonal projection onto $\Alg{G}$
(section~\ref{sec:averages}). When $\mult{p}>1$ the copies are not unique: another choice
$\emb{a}^{p}\mapsto\sum_b\emb{b}^{p}R_{ba}$, $R\in\mathrm O(\mult{p})$, replaces $K^p$ by
$R^{\mathsf T}K^pR$ and changes nothing else. This is Walpole's
construction~\cite{Walpole1984} for fourth-order tensors, used for isotropic sixth-order
tensors by Monchiet and Bonnet~\cite{Monchiet2011a} and Auffray~\cite{Auffray2015}; Glüge
et al.~\cite{Gluge2016} give the spectral decomposition of the isotropic tensor of form~I,
with a harmonic reading of its eigenmodes.

\subsection{The isotropic algebra}\label{sec:isotropic}

For $G=\SOrth$ we first split $\Vsg=\mathcal S\oplus\mathcal R$, with $\mathcal S$ the
totally symmetric tensors (dimension 10, projector $\tP{\mathcal S}\sgrad=\eta_{(ijk)}$)
and $\mathcal R$ its orthogonal complement (dimension 8). This splitting is invariant
under $\GL$~\cite{Auffray2015,Lazar2016}, separates the stretch gradient
($\mathcal S$) from the rotation gradient ($\mathcal R$)~\cite{Auffray2015} and is the
grouping of Mindlin's form~III~\cite{MindlinEshel1968}. Removing the
traces then gives
\begin{equation}
\mathcal S=\Hrm3\oplus\Hone{s},\qquad \mathcal R=\Hrm2\oplus\Hone{r},\qquad
\Vsg=\Hrm3\oplus\Hrm2\oplus2\Hrm1\quad(18=7+5+3+3),\label{eq:iso-decomposition}
\end{equation}
where $\Hrm{\ell}$ has degree $\ell$. $\Hrm3$ is embedded as the totally symmetric
traceless tensors; for vectors $v,w$ and traceless symmetric $X$,
\begin{equation}
\begin{gathered}
(\emb{s}v)_{ijk}=\frac{v_i\delta_{jk}+v_j\delta_{ik}+v_k\delta_{ij}}{\sqrt{15}},\qquad
(\emb{r}w)_{ijk}=\frac{2\delta_{ij}w_k-\delta_{ik}w_j-\delta_{jk}w_i}{\sqrt{12}},\\
(\emb{2}X)_{ijk}=-\frac{\epsilon_{jkl}X_{li}+\epsilon_{ikl}X_{lj}}{\sqrt6},
\end{gathered}\label{eq:iso-embeddings}
\end{equation}
with $\epsilon_{ijk}$ the permutation symbol. The trace $\eta_{ppk}$ of $\emb{s}v$ and
$\emb{r}w$ is a positive multiple of $v$ and $w$, and for a displacement field, with
$\bm\omega=\tfrac12\nabla\times\disp$,
\begin{equation}
\emb{s}^{\mathsf T}\sgrad=\frac{\nabla^2\disp+2\nabla(\nabla\cdot\disp)}{\sqrt{15}},\qquad
\emb{r}^{\mathsf T}\sgrad=\frac{\nabla\times\bm\omega}{\sqrt3},\qquad
\emb{2}^{\mathsf T}\sgrad=\sqrt{\tfrac23}\,\sym\nabla\bm\omega.\label{eq:iso-coordinates}
\end{equation}

$\Alg{\clsiso}$ has dimension 6 (5 symmetric), with the units
\begin{equation}
\tP{3}=\tP{\mathcal S}-\tU{1}{s}{s},\qquad \tP{2}=\tI-\tP{\mathcal S}-\tU{1}{r}{r},\qquad
\tU{1}{a}{b}=\emb{a}\emb{b}^{\mathsf T}\quad(a,b\in\{s,r\}).\label{eq:iso-units}
\end{equation}
An isotropic gradient tensor and its inverse are
\begin{equation}
\tD=m_3\tP{3}+m_2\tP{2}+\sum_{a,b}M_{ab}\tU{1}{a}{b},\qquad
\tD^{-1}=\frac{\tP{3}}{m_3}+\frac{\tP{2}}{m_2}+\sum_{a,b}(M^{-1})_{ab}\tU{1}{a}{b},
\label{eq:iso-tensor}
\end{equation}
with $M$ symmetric $2\times2$; $\tD\succ0$ if and only if $m_3>0$, $m_2>0$ and $M\succ0$.
In the constants $a_1,\dots,a_5$ of Lazar et al.~\cite[eqs.~(40) and (70)]{Lazar2022},
\begin{equation}
\begin{gathered}
m_3=2(a_4+a_5),\qquad m_2=2a_4-a_5,\qquad
M_{sr}=\tfrac{\sqrt5}{3}(a_1+4a_2-2a_3),\\
M_{ss}=\tfrac{10}{3}(a_1+a_2+a_3)+m_3,\qquad
M_{rr}=-\tfrac43a_1+\tfrac83a_2+\tfrac23a_3+m_2.
\end{gathered}\label{eq:iso-lazar}
\end{equation}
The moduli $m_3$, $m_2$, $M_{ss}$ and $M_{rr}$ are $m_{s3}$, $m_{r2}$, $m_{s1}$ and
$m_{r1}$ of Auffray~\cite[table~3]{Auffray2015}, and his coupling $m_{c1}$ is $M_{sr}$ up
to the orientation of $\Hone{r}$. A couple-stress type material, whose energy depends only
on the rotation gradient, has $\tD=\tP{\mathcal R}\tD\tP{\mathcal R}$ and so
$m_3=M_{ss}=M_{sr}=0$: these copies leave it uncoupled. Lazar~\cite{Lazar2016} uses
another pair; appendix~\ref{app:lazar} gives the conversion.

\section{Intrinsic basis of the cubic algebra}\label{sec:cubic}

\subsection{Branching and embeddings}\label{sec:cubic-embeddings}

Let $\bm e_1,\bm e_2,\bm e_3$ be unit vectors along the cube axes and
$\tOct=\sum_{s}\bm e_s\ot\bm e_s\ot\bm e_s\ot\bm e_s$ the cubic structure tensor (the
$\delta_{ijkl}$ of Lazar et al.~\cite[eq.~(37)]{Lazar2022}). Under the octahedral group
the isotropic components branch without multiplicity,
$\Hrm3\to\irr{A_2}\oplus\irr{T_1}\oplus\irr{T_2}$, $\Hrm2\to\irr{E}\oplus\irr{T_2}$ and
$\Hrm1\to\irr{T_1}$ (Mulliken labels), so that
\begin{equation}
\Vsg=3\irr{T_1}\oplus2\irr{T_2}\oplus\irr{E}\oplus\irr{A_2}\qquad(18=3\cdot3+2\cdot3+2+1).
\label{eq:cub-decomposition}
\end{equation}
As model spaces we take the vectors for $\irr{T_1}$, the symmetric matrices with
$\tOct:X=0$ (zero diagonal in the cube axes) for $\irr{T_2}$, the traceless symmetric
matrices with $\tOct:X=X$ (diagonal) for $\irr{E}$, and a line for $\irr{A_2}$. The copies
are ordered by decreasing degree and $\mathcal S$ before $\mathcal R$:
$\irr{T_1}=(h,s,r)$ and $\irr{T_2}=(3,2)$. The copies $s$ and $r$ of $\irr{T_1}$ are
$\emb{s}$ and $\emb{r}$ of~\eqref{eq:iso-embeddings}; copy~2 of $\irr{T_2}$ and the
$\irr{E}$ component are $\emb{2}$ restricted to their model spaces. The copies inside
$\Hrm3$ are
\begin{equation}
\begin{gathered}
(\emb{h}v)_{ijk}=\sqrt{\tfrac52}\Big[\tOct_{ijkl}v_l
-\tfrac15(v_i\delta_{jk}+v_j\delta_{ik}+v_k\delta_{ij})\Big],\\
(\emb{3}X)_{ijk}=\sqrt3\,\big[\epsilon_{ilp}X_{pq}\tOct_{lqjk}\big]_{(ijk)},
\end{gathered}
\label{eq:cub-embeddings}
\end{equation}
where $(\cdot)_{(ijk)}$ is total symmetrisation, and $\irr{A_2}$ is the line of $\Hrm3$
annihilated by $h\mapsto \tOct_{ijpq}h_{pqk}$, spanned by the symmetrised
$\bm e_1\ot\bm e_2\ot\bm e_3$. Thus $\emb{h}v$ is a positive multiple of the harmonic part of $\tOct\cdot v$, and $\emb{3}$ is oriented so that the separable tensor
$\tOct_{ijkl}\delta_{mn}$ (the $a_8$ term of~\cite[eq.~(40)]{Lazar2022}) has the
positive coupling $N_{32}=\sqrt2/3$ between the two copies of $\irr{T_2}$.

\subsection{The algebra and the cubic tensor}\label{sec:cubic-algebra}

$\Alg{\clscub}$ has dimension $9+4+1+1=15$, with 11 symmetric elements
(table~\ref{tab:algebras}); its units $\tU{T_1}{a}{b}$ ($a,b\in\{h,s,r\}$),
$\tU{T_2}{a}{b}$ ($a,b\in\{3,2\}$), $\tP{E}$ and $\tP{A_2}$ obey~\eqref{eq:table}, and the
dimensions agree with those of the commutant computed directly from the group action. The
units are tied to the isotropic ones by
\begin{equation}
\tP{3}=\tU{T_1}{h}{h}+\tU{T_2}{3}{3}+\tP{A_2},\qquad
\tP{2}=\tU{T_2}{2}{2}+\tP{E},\qquad
\tU{1}{a}{b}=\tU{T_1}{a}{b}\quad(a,b\in\{s,r\}),\label{eq:cub-iso-units}
\end{equation}
with $\tP{E}=\emb{2}\,\tOct\,\emb{2}^{\mathsf T}$, where $\tOct$ projects a traceless
symmetric matrix onto its diagonal part. Hence $\tU{T_2}{2}{2}=\tP{2}-\tP{E}$ and
$\tP{A_2}=\tP{3}-\tU{T_1}{h}{h}-\tU{T_2}{3}{3}$, and every unit is a composition of maps
built from $\kron$, $\epsilon$ and $\tOct$; their expansions in components are in the
supplementary material. A cubic gradient tensor is
\begin{equation}
\tD=\sum_{a,b}M_{ab}\,\tU{T_1}{a}{b}+\sum_{a,b}N_{ab}\,\tU{T_2}{a}{b}
+m_E\,\tP{E}+m_A\,\tP{A_2},\label{eq:cub-tensor}
\end{equation}
with $M=M^{\mathsf T}\in\mathbb R^{3\times3}$ and $N=N^{\mathsf T}\in\mathbb R^{2\times2}$,
that is, $6+3+1+1=11$ block moduli. By~\eqref{eq:block-rules} its eigenvalues are those of
$M$ and of $N$, each threefold, $m_E$, twofold, and $m_A$, simple: the seven distinct
eigenvalues (one simple, one double, five threefold) that Lazar et
al.~\cite{Lazar2022} obtain by reordering the $18\times18$ matrix into diagonal blocks of
sizes five and three. Each of their blocks of size five joins the $\irr{T_1}$ and
$\irr{T_2}$ parts, which the basis separates into blocks of sizes three and two. Products, inverses and
positive definiteness follow from~\eqref{eq:block-rules}; section~\ref{sec:compliance}
states them. By~\eqref{eq:cub-iso-units}, $\tD$ is isotropic if and only if
\begin{equation}
M_{hh}=N_{33}=m_A,\qquad N_{22}=m_E,\qquad M_{hs}=M_{hr}=N_{32}=0,
\label{eq:cub-isotropy}
\end{equation}
and then $m_3=m_A$, $m_2=m_E$ and the $(s,r)$ part of $M$ is the isotropic block. These
six linear conditions match the six conditions $a_6=\dots=a_{11}=0$
of~\cite[eq.~(81)]{Lazar2022}.

\subsection{The constants of Lazar et al.}\label{sec:cubic-lazar}

Lazar et al.~\cite[eq.~(40)]{Lazar2022} write the cubic tensor with eleven constants:
$a_1,\dots,a_5$ multiply Mindlin's isotropic tensors and $a_6,\dots,a_{11}$ tensors built
from $\kron$, $\tOct$ and $\sum_s\bm e_s^{\ot6}$. Their tensor lies in $\Alg{\clscub}$, and
its block moduli are obtained from~\eqref{eq:block-rules}. The moduli $m_E$ and $m_A$ are
the isotropic $m_2$ and $m_3$ of~\eqref{eq:iso-lazar} and do not depend on the anisotropic
constants; with $m_3$, $m_2$ and $M^{\clsiso}_{ab}$ the isotropic moduli
of~\eqref{eq:iso-lazar}, the other blocks are
\begin{equation}
\begin{aligned}
M_{hh}&=m_3+\tfrac{12}{5}(a_6+a_7)+\tfrac35a_8+\tfrac25a_{11},&
N_{33}&=m_3+\tfrac43(a_6+a_7)+\tfrac13a_8,\\
M_{ss}&=M^{\clsiso}_{ss}+\tfrac{44}{15}(a_6+a_7)+\tfrac{11}{15}a_8+2a_9+4a_{10}+\tfrac35a_{11},&
N_{22}&=m_2+\tfrac23a_6-\tfrac43a_7+\tfrac23a_8,\\
M_{rr}&=M^{\clsiso}_{rr}+\tfrac23a_6-\tfrac43a_7+\tfrac23a_8,&
N_{32}&=\tfrac{\sqrt2}{3}(-2a_6+a_7+a_8),\\
M_{sr}&=M^{\clsiso}_{sr}+\sqrt5\big[-\tfrac4{15}a_6+\tfrac2{15}(a_7+a_8)+\tfrac25(a_9-a_{10})\big],\\
M_{hs}&=\sqrt6\big[\tfrac8{15}(a_6+a_7)+\tfrac2{15}a_8+\tfrac13a_9+\tfrac23a_{10}+\tfrac15a_{11}\big],\\
M_{hr}&=\tfrac{\sqrt{30}}{15}(2a_6-a_7-a_8+2a_9-2a_{10}).
\end{aligned}\label{eq:cub-lazar}
\end{equation}
The map from $(a_1,\dots,a_{11})$ to the eleven block moduli is invertible; its inverse,
and the conditions of section~\ref{sec:compliance} in terms of the $a_i$, are given there.
The atomistic data of section~\ref{sec:compliance} enter through~\eqref{eq:cub-lazar}.

\begin{table}[pos=tbp]
\caption{Structure of the isotropic, cubic and hexagonal algebras on the strain-gradient
space $\Vsg$. For each irreducible representation $p$ of the rotation group ($\SOrth$,
$O$ or $D_6$): its dimension $\dimirr{p}$; the isotropic components that contain its
copies, in the order of the blocks; its multiplicity $\mult{p}$; and its contributions to
the dimension of the algebra ($\mult{p}^2$) and to the number of constants of the tensors
with major symmetry ($\tfrac12\mult{p}(\mult{p}+1)$). Hexagonal labels refer to $\bm e_1$
along an $a$ axis (section~\ref{sec:hexagonal}).}\label{tab:algebras}
\centering\small
\input{algebra_structure}
\end{table}

\section{Intrinsic basis of the hexagonal algebra and the extra modulus}\label{sec:hexagonal}

\subsection{Structure tensors and embeddings}\label{sec:hex-embeddings}

Let $\bm n$ be a unit vector along the sixfold axis, $\bm a$ a unit vector along one of
the twofold axes normal to it, and $\zhex=\bm a+\iu\,\bm n\times\bm a$. Since
$\zhex\cdot\zhex=0$, the tensors
\begin{equation}
\bm h=\mathrm{Re}\,\zhex^{\ot3},\qquad \bm k=\mathrm{Im}\,\zhex^{\ot3},\qquad
\tHsix=\mathrm{Re}\,\zhex^{\ot6}\label{eq:hex-structure}
\end{equation}
are totally symmetric and traceless; $|\bm h|=|\bm k|=2$, and $\tHsix$ is the
sixth-order hexagonal structure tensor. A rotation by $\varphi$ about $\bm n$ maps $\zhex$
to $e^{-\iu\varphi}\zhex$: by $\pi/3$ it reverses $\bm h$ and $\bm k$, by $\pi/6$ it
maps $\bm h$ to $\bm k$ and $\bm k$ to $-\bm h$. Under $D_6$ the isotropic components
branch without multiplicity, $\Hrm3\to\irr{A_2}\oplus\irr{B_1}\oplus\irr{B_2}\oplus
\irr{E_1}\oplus\irr{E_2}$, $\Hrm2\to\irr{A_1}\oplus\irr{E_1}\oplus\irr{E_2}$ and
$\Hrm1\to\irr{A_2}\oplus\irr{E_1}$, with $\irr{B_1}$ even under the half-turn about
$\bm a$; hence
\begin{equation}
\Vsg=\irr{A_1}\oplus3\irr{A_2}\oplus\irr{B_1}\oplus\irr{B_2}\oplus4\irr{E_1}\oplus
2\irr{E_2}\qquad(18=1+3+1+1+4\cdot2+2\cdot2).\label{eq:hex-decomposition}
\end{equation}
The model spaces are lines for the one-dimensional representations, the vectors
$\bm v\perp\bm n$ for $\irr{E_1}$ and the traceless symmetric matrices $X$ with
$X\bm n=\bm0$ for $\irr{E_2}$. The copies are ordered as in section~\ref{sec:cubic}:
$\irr{A_2}=(3,s,r)$, $\irr{E_1}=(3,2,s,r)$ and $\irr{E_2}=(3,2)$. The copies $s$ and $r$
are $\emb{s}$ and $\emb{r}$ of~\eqref{eq:iso-embeddings} applied to $\bm n$ and to
$\bm v$, and the copies~2 are $\emb{2}$ applied to $\sqrt{3/2}\,(\bm n\ot\bm
n-\tfrac13\kron)$ ($\irr{A_1}$), to $\sqrt2\,\sym(\bm n\ot(\bm n\times\bm v))$
($\irr{E_1}$) and to $X$ ($\irr{E_2}$). With $\harm$ the harmonic (totally symmetric
traceless) part, the copies inside $\Hrm3$ are
\begin{equation}
\begin{gathered}
\emb{3}^{A_2}=\sqrt{\tfrac52}\,\harm(\bm n\ot\bm n\ot\bm n),\qquad
\emb{}^{B_1}=\tfrac12\bm h,\qquad \emb{}^{B_2}=\tfrac12\bm k,\\
\emb{3}^{E_1}\bm v=\tfrac{\sqrt{15}}{2}\,\harm(\bm n\ot\bm n\ot\bm v),\qquad
(\emb{3}^{E_2}X)_{ijk}=\sqrt3\,\big[n_i\,\epsilon_{jpq}n_pX_{qk}\big]_{(ijk)}.
\end{gathered}\label{eq:hex-embeddings}
\end{equation}
$\emb{3}^{A_2}$ is $\emb{h}$ of~\eqref{eq:cub-embeddings} with $\tOct$ replaced by
$\bm n^{\ot4}$ and $\bm v=\bm n$. The 18 embedded basis vectors have the identity as
Gram matrix.

\subsection{The algebra and the hexagonal tensor}\label{sec:hex-algebra}

$\Alg{\clshex}$ has dimension $1+9+1+1+16+4=32$, with 22 symmetric elements
(table~\ref{tab:algebras}); both numbers equal the dimensions of the commutant of $D_6$
computed directly. The
isotropic units split as
\begin{equation}
\tP{3}=\tU{A_2}{3}{3}+\tP{B_1}+\tP{B_2}+\tU{E_1}{3}{3}+\tU{E_2}{3}{3},\quad
\tP{2}=\tP{A_1}+\tU{E_1}{2}{2}+\tU{E_2}{2}{2},\quad
\tU{1}{a}{b}=\tU{A_2}{a}{b}+\tU{E_1}{a}{b},\label{eq:hex-iso-units}
\end{equation}
with $a,b\in\{s,r\}$ in the last one. A hexagonal gradient tensor is
\begin{equation}
\tD=\sum_{a,b}\hexA_{ab}\tU{A_2}{a}{b}+\sum_{a,b}\hexF_{ab}\tU{E_1}{a}{b}
+\sum_{a,b}\hexG_{ab}\tU{E_2}{a}{b}+m_{A_1}\tP{A_1}+m_{B_1}\tP{B_1}+m_{B_2}\tP{B_2},
\label{eq:hex-tensor}
\end{equation}
with symmetric blocks $\hexA$ ($3\times3$), $\hexF$ ($4\times4$) and $\hexG$ ($2\times2$):
$6+10+3+3=22$ block moduli. Its eigenvalues are those of $\hexA$, $m_{A_1}$, $m_{B_1}$
and $m_{B_2}$, simple, and those of $\hexF$ and $\hexG$, double. By~\eqref{eq:hex-iso-units},
$\tD$ is isotropic if and only if
\begin{equation}
\begin{gathered}
m_{B_1}=m_{B_2}=\hexA_{33}=\hexF_{33}=\hexG_{33},\qquad m_{A_1}=\hexF_{22}=\hexG_{22},
\qquad \hexA_{ab}=\hexF_{ab}\quad(a,b\in\{s,r\}),\\
\hexA_{3s}=\hexA_{3r}=\hexF_{32}=\hexF_{3s}=\hexF_{3r}=\hexF_{2s}=\hexF_{2r}=\hexG_{32}=0,
\end{gathered}\label{eq:hex-isotropy}
\end{equation}
seventeen conditions, and then $m_3=m_{B_1}$, $m_2=m_{A_1}$ and the $(s,r)$ blocks are
$M$.

\subsection{Transverse isotropy and the extra modulus}\label{sec:hex-extra}

All the units except $\tP{B_1}$ and $\tP{B_2}$ commute with every rotation about
$\bm n$, and so does their sum, which is the projector onto the part of $\Hrm3$ of azimuthal order $m=\pm3$ (the
span of $\bm h$ and $\bm k$). The transversely isotropic tensors therefore form the
subalgebra in which $\tP{B_1}$ and $\tP{B_2}$ enter only through their sum, of dimension
31 (21 symmetric), spanned by the 31 tensors of Monchiet and
Bonnet~\cite[eqs.~(32)--(33)]{Monchiet2013}; a hexagonal tensor is transversely
isotropic if and only if $\Dhex:=m_{B_1}-m_{B_2}=0$. Since $\tP{B_1}=\tfrac14\bm h\ot\bm
h$, $\tP{B_2}=\tfrac14\bm k\ot\bm k$ and $\mathrm{Re}\,u\,\mathrm{Re}\,w-\mathrm{Im}\,u\,
\mathrm{Im}\,w=\mathrm{Re}\,(uw)$,
\begin{equation}
\tP{B_1}-\tP{B_2}=\tfrac14\,\tHsix,\qquad
\Dhex=m_{B_1}-m_{B_2}=\tfrac14\,\langle\tD,\tHsix\rangle=\tfrac14\,D_{ijmkln}\,
(\tHsix)_{ijmkln}.\label{eq:hex-extra}
\end{equation}
The single modulus that distinguishes the class $\clshex$ from transverse isotropy is
thus the splitting of the $m=\pm3$ part of $\Hrm3$ into its $\irr{B_1}$ and $\irr{B_2}$
parts, and, up to a factor, the component of $\tD$ along the hexagonal structure tensor:
since $|\tHsix|^2=32$, the orthogonal projection of $\tD$ on $\tHsix$ is
$\tfrac18\Dhex\,\tHsix$. A harmonic tensor
invariant under the rotation by $\pi/3$ about $\bm n$ has only components with
$m\equiv0\pmod 6$, so the $D_6$-invariant harmonic tensors of degree below six are
axially symmetric. This is why a classical elasticity tensor of class $\clshex$, whose
harmonic components have degree at most four, is transversely isotropic, whereas the
symmetric gradient tensors, whose space contains the degree-six component of
$\Hrm3\ot\Hrm3$, need not be. The rotation by $\pi/6$ about $\bm n$ exchanges $\tP{B_1}$ and
$\tP{B_2}$ and reverses $\Dhex$, so the sign of $\Dhex$ refers to the family of twofold
axes that contains $\bm a$. A separable tensor $\tC\ot\bm L$, with $\tC$ hexagonal and
$\bm L=l_1\bm n\ot\bm n+l_2(\kron-\bm n\ot\bm n)$, has $\Dhex=0$. Auffray et
al.~\cite[section~3.3.5]{Auffray2013} write the hexagonal matrix as a transversely
isotropic matrix plus $\eta_{\mathrm A}A_c$, with a fixed matrix $A_c$; in our basis
$\eta_{\mathrm A}A_c$ is $4\eta_{\mathrm A}\,\tP{B_1}$, so that $\Dhex=4\eta_{\mathrm A}$ when both are referred to the same
$\bm e_1$.

\subsection{The matrix components of Auffray et al.}\label{sec:hex-auffray}

Auffray et al.~\cite[eq.~(10) and table~2]{Auffray2013} use an orthonormal basis of
$\Vsg$ in which the hexagonal matrix is block diagonal, with blocks
$A^{(11)}+\eta_{\mathrm A}A_c$, $A^{(11)}$, $H^{(6)}+f(J^{(4)})$ and $J^{(4)}$ of sizes 5, 5, 5 and
3~\cite[section~3.3.5]{Auffray2013}. With $\bm e_3=\bm n$ and $\bm e_1=\bm a$, the
change to our basis is orthogonal and respects these blocks:
\begin{equation}
\begin{gathered}
A^{(11)}+\eta_{\mathrm A}A_c=W_1\begin{pmatrix}m_{B_1}&\\&\hexF\end{pmatrix}W_1^{\mathsf T},\qquad
A^{(11)}=W_1'\begin{pmatrix}m_{B_2}&\\&\hexF\end{pmatrix}W_1'^{\mathsf T},\\
H^{(6)}+f(J^{(4)})=W_3\begin{pmatrix}\hexA&\\&\hexG\end{pmatrix}W_3^{\mathsf T},\qquad
J^{(4)}=W_4\begin{pmatrix}m_{A_1}&\\&\hexG\end{pmatrix}W_4^{\mathsf T},\\[2pt]
W_1=\frac1{30}\begin{pmatrix}
15&-3\sqrt{15}&0&6\sqrt{15}&0\\
-15&-\sqrt{15}&10\sqrt3&2\sqrt{15}&10\sqrt3\\
-15\sqrt2&-\sqrt{30}&-5\sqrt6&2\sqrt{30}&-5\sqrt6\\
0&4\sqrt{15}&-10\sqrt3&2\sqrt{15}&10\sqrt3\\
0&4\sqrt{30}&5\sqrt6&2\sqrt{30}&-5\sqrt6\end{pmatrix},\quad
W_4=\frac16\begin{pmatrix}0&2\sqrt3&-2\sqrt6\\-3\sqrt2&2\sqrt3&\sqrt6\\
3\sqrt2&2\sqrt3&\sqrt6\end{pmatrix},\\[2pt]
W_3=\frac1{30}\begin{pmatrix}
6\sqrt{10}&6\sqrt{15}&0&0&0\\
-3\sqrt{10}&2\sqrt{15}&10\sqrt3&-5\sqrt6&10\sqrt3\\
-6\sqrt5&2\sqrt{30}&-5\sqrt6&-10\sqrt3&-5\sqrt6\\
-3\sqrt{10}&2\sqrt{15}&10\sqrt3&5\sqrt6&-10\sqrt3\\
-6\sqrt5&2\sqrt{30}&-5\sqrt6&10\sqrt3&5\sqrt6\end{pmatrix},
\end{gathered}\label{eq:hex-auffray}
\end{equation}
where $W_1'$ is $W_1$ with the sign of its first column reversed, which changes only the
orientation of $\emb{}^{B_2}$. The rows follow their basis
vectors (111, 221, 122, 331, 133), (222, 112, 121, 332, 233), (333, 113, 131, 223, 232)
and (123, 132, 231), and the columns our copies, $\irr{B_1}$ or $\irr{B_2}$ before
$\irr{E_1}$ and $\irr{A_2}$ or $\irr{A_1}$ before $\irr{E_2}$. With $\bm b=\bm n\times\bm a$,
the $\irr{E_1}$ columns take $\bm v=\bm a$ in the first block and $\bm v=\bm b$ in the
second, and the $\irr{E_2}$ columns take $X=\sqrt2\,\sym(\bm a\ot\bm b)$ in the third and
$X=(\bm a\ot\bm a-\bm b\ot\bm b)/\sqrt2$ in the fourth. Since the $W$ are
orthogonal, \eqref{eq:hex-auffray} gives both directions, for instance
$\mathrm{diag}(m_{B_2},\hexF)=W_1'^{\mathsf T}A^{(11)}W_1'$, and in particular $m_{A_1}=j_{22}-j_{23}$ and $\hexG_{32}=\tfrac{\sqrt2}{3}(-j_{11}-j_{12}+j_{22}
+j_{23})$. Each of their blocks thus mixes two of our irreducible components:
$\irr{B_1}$ or $\irr{B_2}$ with $\irr{E_1}$ in the first two, $\irr{A_2}$ or
$\irr{A_1}$ with $\irr{E_2}$ in the last two.

\section{Closed-form gradient compliance}\label{sec:compliance}

\subsection{Inverse and conditions}\label{sec:compliance-inverse}

By~\eqref{eq:block-rules}, the compliances of the cubic tensor~\eqref{eq:cub-tensor} and
of the hexagonal tensor~\eqref{eq:hex-tensor} are
\begin{equation}
\begin{aligned}
\tD^{-1}&=\sum_{a,b}(M^{-1})_{ab}\tU{T_1}{a}{b}+\sum_{a,b}(N^{-1})_{ab}\tU{T_2}{a}{b}
+\frac{\tP{E}}{m_E}+\frac{\tP{A_2}}{m_A},\\
\tD^{-1}&=\sum_{a,b}(\hexA^{-1})_{ab}\tU{A_2}{a}{b}+\sum_{a,b}(\hexF^{-1})_{ab}\tU{E_1}{a}{b}
+\sum_{a,b}(\hexG^{-1})_{ab}\tU{E_2}{a}{b}
+\frac{\tP{A_1}}{m_{A_1}}+\frac{\tP{B_1}}{m_{B_1}}+\frac{\tP{B_2}}{m_{B_2}}.
\end{aligned}\label{eq:compliance}
\end{equation}
The blocks have size at most three (cubic) and four (hexagonal), so their inverses are
explicit by cofactors, and the units are explicit in components (supplementary
material, section~S1); the same formulas invert the non-symmetric elements of the
algebras, whose blocks are not symmetric. A cubic tensor is invertible if and only if
$\det M\,\det N\,m_Em_A\ne0$, and positive definite if and only if, by Sylvester's
criterion on the blocks,
\begin{equation}
M_{hh}>0,\quad M_{hh}M_{ss}-M_{hs}^2>0,\quad \det M>0,\quad N_{33}>0,\quad \det N>0,
\quad m_E>0,\quad m_A>0.\label{eq:cub-pd}
\end{equation}
These are seven conditions, as many as Lazar et al.\ obtain from Sylvester's criterion
on the $5\times5$ and $3\times3$ diagonal blocks of their matrix~\cite[appendix~A,
eqs.~(A.7)]{Lazar2022}, but each involves a single irreducible block and has degree at
most three in the moduli. A hexagonal tensor is invertible if and only if $\det\hexA\,\det\hexF\,\det\hexG\,
m_{A_1}m_{B_1}m_{B_2}\ne0$ and positive definite if and only if $\hexA$, $\hexF$ and
$\hexG$ are positive definite and $m_{A_1}$, $m_{B_1}$, $m_{B_2}$ are positive; in the
components of Auffray et al., \eqref{eq:hex-auffray} turns these into the positive
definiteness of $A^{(11)}$, $A^{(11)}+\eta_{\mathrm A}A_c$, $H^{(6)}+f(J^{(4)})$ and $J^{(4)}$.
By~\eqref{eq:hex-extra},
\begin{equation}
\frac{\tP{B_1}}{m_{B_1}}+\frac{\tP{B_2}}{m_{B_2}}
=\frac12\Big(\frac1{m_{B_1}}+\frac1{m_{B_2}}\Big)(\tP{B_1}+\tP{B_2})
+\frac18\Big(\frac1{m_{B_1}}-\frac1{m_{B_2}}\Big)\tHsix,\label{eq:hex-extra-compliance}
\end{equation}
so the extra modulus of $\tD^{-1}$ is $-\Dhex/(m_{B_1}m_{B_2})$: the compliance of a
transversely isotropic tensor is transversely isotropic, and that of an isotropic tensor
is the isotropic inverse~\eqref{eq:iso-tensor}.

\subsection{Back to the constants of Lazar et al.}\label{sec:compliance-lazar}

The map~\eqref{eq:cub-lazar} is inverted in closed form. With
$u=M_{hh}-m_A$, $x=N_{33}-m_A$, $y=N_{22}-m_E$ and
$X=M_{ss}-\sqrt6M_{hs}+\tfrac32M_{hh}-\tfrac52N_{33}$, $Y=M_{rr}-N_{22}$,
$Z=M_{sr}-\tfrac{\sqrt6}{2}M_{hr}-\tfrac{\sqrt{10}}{2}N_{32}$,
\begin{equation}
\begin{gathered}
a_1=\tfrac1{15}\big(2X-5Y+\sqrt5Z\big),\quad
a_2=\tfrac1{30}\big(X+5Y+2\sqrt5Z\big),\quad
a_3=\tfrac1{30}\big(4X+5Y-4\sqrt5Z\big),\\
a_4=\tfrac16(m_A+2m_E),\quad a_5=\tfrac13(m_A-m_E),\quad
a_6=\tfrac16\big(2x+y-2\sqrt2N_{32}\big),\\
a_7=\tfrac16\big(2x-2y+\sqrt2N_{32}\big),\quad
a_8=\tfrac13\big(x+2y+2\sqrt2N_{32}\big),\quad
a_{11}=\tfrac12(5u-9x),\\
a_9=\tfrac12(x-u)+\tfrac16\big(\sqrt6M_{hs}+\sqrt{30}M_{hr}+3\sqrt2N_{32}\big),\\
a_{10}=\tfrac12(x-u)+\tfrac1{12}\big(2\sqrt6M_{hs}-\sqrt{30}M_{hr}-3\sqrt2N_{32}\big).
\end{gathered}\label{eq:cub-lazar-inverse}
\end{equation}
For an isotropic tensor $u=x=y=0$, $X=M_{ss}-m_3$, $Y=M_{rr}-m_2$ and $Z=M_{sr}$, and
the first line, with $a_4$ and $a_5$, inverts~\eqref{eq:iso-lazar}. The conditions
\eqref{eq:cub-pd} in the constants $a_i$ follow by inserting~\eqref{eq:cub-lazar}.

\subsection{Atomistic data}\label{sec:compliance-data}

Table~\ref{tab:block-moduli} gives the block moduli of five cubic data sets: the MEAM
constants of Al, Cu, Fe and W of Lazar et al.~\cite[table~1]{Lazar2022}, converted to
$a_1,\dots,a_{11}$ by their eq.~(60), and the EAM constants of Cu of Po et
al.~\cite[table~1]{Po2019}. Po et al.\ print the Al and Cu MEAM constants too, with five
significant figures instead of five decimals; for the entries below $0.1$ in magnitude
their digits are finer, and we take each entry from the finer of the two printings. Shodja et al.~\cite{Shodja2018} evaluated cubic constants by
density functional theory under a central-force assumption, which imposes equalities
among them; we
use the complete tensors. The data of Al and of both potentials for Cu are positive
definite. Those of Fe and W are not: Lazar et al.\ find one negative eigenvalue, threefold, a
root of the quadratic factor~\cite[eqs.~(130)--(133) and table~5]{Lazar2022}. In the basis
that factor is the $\irr{T_2}$ block, the only one that fails, with $\det N=\NumFeDetNT$ and $\NumWDetNT$~(eV/\AA)$^2$. Its
entry $N_{33}$, the $\irr{T_2}$ part of $\Hrm3$, is negative ($\NumFeNthreethreeT$ and
$\NumWNthreethreeT$~eV/\AA), but for W it nearly vanishes, and the indefiniteness comes
from the coupling $N_{32}=\NumWNthreetwoT$~eV/\AA. With~\eqref{eq:cub-lazar}, the two
conditions $N_{33}>0$ and $\det N>0$ are explicit in the constants $a_i$, and a fit of a
potential could impose them. Indefinite gradient tensors
are common among interatomic potentials~\cite{Admal2017,Lazar2022}; Admal et
al.~\cite[eq.~(3.11)]{Admal2017} propose instead the positive definiteness of
$\tfrac15R^2\,\tC\ot\kron+\tD$ for a sphere of radius $R$. Both terms are cubic, so the
condition is a generalised eigenvalue problem per block; it holds for
$R>\NumFeNeighborhoodRadius$~\AA\ (Fe) and $R>\NumWNeighborhoodRadius$~\AA\ (W), that is
$\NumFeNeighborhoodRatio\,b$ and $\NumWNeighborhoodRatio\,b$ with $b$ the
nearest-neighbour distance, consistent with their finding that it holds beyond $1.2\,b$ for
all the potentials they examine, Fe and W included~\cite[section~3]{Admal2017}. The Cu MEAM tensor
is nearly singular, with condition number $\NumCuMEAMCond$, because its $m_E$ is close to
zero; the rounding of the published constants fixes $m_E$ only to within
$\pm\NumCuMEAMmERoundingPercent\,\%$. These are properties of the potentials, not of the
method. The compliance blocks of the three positive definite sets are in the
supplementary material (table~S3).

\subsection{Verification}\label{sec:compliance-verification}

The results were checked by routes other than those that produced them. Blocks, inverses
and averages are compared with the same operations on the $18\times18$ matrix of the
tensor restricted to $\Vsg$; the multiplication tables, the inverses of general elements
and the closed averages are also verified in exact arithmetic; the formulas printed in
sections~\ref{sec:setting}--\ref{sec:averages} and appendix~\ref{app:lazar} are restated
in the tests as printed and checked against the code. Published results enter after reconstruction: the constants
and eigenvalues of Lazar et al.~\cite{Lazar2022}, the matrices of Auffray
et al.~\cite{Auffray2013} and the tensors of Monchiet and Bonnet~\cite{Monchiet2011a,Monchiet2013}
are reproduced (except the eigenvalue $\lambda_7^D$ of Cu in~\cite[table~5]{Lazar2022},
printed as a repetition of $\lambda_4^D$; the $\irr{T_1}$ block gives
$\NumCuMEAMLambdaSeven$~eV/\AA), and where a printed formula does not hold, both the printed and the
corrected form are checked. Three such formulas enter this paper: the isotropic inverse of
Monchiet and Bonnet~\cite[eq.~(2.13)]{Monchiet2011a}, whose terms in their $J_3$ and $J_4$
need the opposite signs; the average over the sphere of the sixth-order Eshelby tensor of
Ma et al.~\cite[eq.~(3.18)]{Ma2018} (appendix~\ref{app:eshelby}); and, in the scheme of
Solyaev~\cite{Solyaev2023}, the deviatoric term of the fourth-order average in his eq.~(20),
half that of Ma et al., his eq.~(21), which inherits their eq.~(3.18), and his eq.~(22), printed without the
squares (section~\ref{sec:application}). We distinguish three levels: verification, against exact
mathematics or an independent route of our own; cross-comparison, against a
reconstructed published result; and consistency, through limit cases, invariances and
inequalities on samples. No comparison is with measurements: the inputs are published
atomistic constants and, in the hexagonal class, tensors of granular micromechanics. The tables, figures and
computed numbers are written by scripts of the public code (Data accessibility), whose
tests compare each file with its generator; section~S2 of the supplementary material
maps each claim to its key tests, their reference and their level.

\begin{table}[pos=tbp]
\caption{Block moduli~\eqref{eq:cub-tensor} of the five cubic data sets (eV/\AA), from
the MEAM constants of Lazar et al.~\cite[table~1]{Lazar2022} and the EAM constants of
Po et al.~\cite[table~1]{Po2019} by~\eqref{eq:cub-lazar}; for Al and Cu MEAM each entry
is taken from the finer of the two printings. Last rows: positive
definiteness, with the failing block, and condition number of $\tD$ (largest eigenvalue over
smallest eigenvalue modulus). The last digit shown is uncertain by at most five units
given the precision of the published constants; $\dagger$ marks the entries whose
relative uncertainty exceeds 1\,\%, here the $m_E$ of Cu MEAM ($\pm\NumCuMEAMmERoundingPercent\,\%$)
and the condition number it sets. The Cu EAM constants obey the relations
$D_{2,4}=D_{3,5}$ and $D_{2,5}=D_{16,17}$ that Admal et al.~\cite{Admal2017} find for
pair functionals, that is, $a_3=4a_2$ and $a_1=a_5$.}\label{tab:block-moduli}
\centering\small
\input{block_moduli}
\end{table}

\section{Orientation averages}\label{sec:averages}

\subsection{The rule of weighted means}\label{sec:averages-rule}

In an untextured polycrystal the grain orientations are uniformly distributed, and the
average of an element of a crystal algebra over the rotations,
$\avg{\tK}=\int_{\SOrth}(Q\star)\,\tK\,(Q\star)^{\mathsf T}\,\dd Q$, is isotropic. The
average is linear, idempotent, self-adjoint for $\langle\cdot,\cdot\rangle$ and fixes
$\Alg{\clsiso}$, so it is the orthogonal projection onto $\Alg{\clsiso}$, and the last
formula of~\eqref{eq:block-rules} gives its moduli. By~\eqref{eq:cub-iso-units}
and~\eqref{eq:hex-iso-units}, each isotropic unit is a sum of crystal units; hence the
isotropic moduli of $\avg{\tK}$ are means of the crystal block entries weighted by the
dimensions of the irreducible components,
\begin{equation}
m_\ell=\sum_p\frac{\dimirr{p}}{2\ell+1}\,K^p_{c_\ell c_\ell}\quad(\ell=3,2),\qquad
M_{ab}=\sum_p\frac{\dimirr{p}}{3}\,K^p_{ab}\quad(a,b\in\{s,r\}),\label{eq:weighted-means}
\end{equation}
where $c_\ell$ is the copy of $\Hrm{\ell}$ in $p$ and the sums run over the $p$ that
contain one. The weights add up to one, the couplings between copies of different
$\Hrm{\ell}$ do not survive, and the rule holds for the non-symmetric elements too. For
the Voigt estimate $\Vg{\tD}=\avg{\tD}$ of a cubic tensor,
\begin{equation}
\Vg{m_3}=\tfrac17\big(3M_{hh}+3N_{33}+m_A\big),\qquad
\Vg{m_2}=\tfrac15\big(3N_{22}+2m_E\big),\qquad \Vg{M}=M_{(sr)},\label{eq:cub-voigt}
\end{equation}
with $M_{(sr)}$ the $(s,r)$ part of $M$. The Reuss estimate
$\Rs{\tD}=\avg{\tD^{-1}}^{-1}$ applies the rule to the inverse blocks and inverts each
isotropic block; with $\bm m_h=(M_{hs},M_{hr})^{\mathsf T}$,
\begin{equation}
\begin{gathered}
\frac1{\Rs{m_3}}=\frac17\Big(3(M^{-1})_{hh}+3(N^{-1})_{33}+\frac1{m_A}\Big),\qquad
\frac1{\Rs{m_2}}=\frac15\Big(3(N^{-1})_{22}+\frac2{m_E}\Big),\\
\Rs{M}=M_{(sr)}-\frac{\bm m_h\bm m_h^{\mathsf T}}{M_{hh}},
\end{gathered}\label{eq:cub-reuss}
\end{equation}
the last being the Schur complement of $M_{hh}$, the inverse of $(M^{-1})_{(sr)}$. For a
hexagonal tensor,
\begin{equation}
\begin{gathered}
\Vg{m_3}=\tfrac17\big(\hexA_{33}+m_{B_1}+m_{B_2}+2\hexF_{33}+2\hexG_{33}\big),\qquad
\Vg{m_2}=\tfrac15\big(m_{A_1}+2\hexF_{22}+2\hexG_{22}\big),\\
\Vg{M}=\tfrac13\big(\hexA_{(sr)}+2\hexF_{(sr)}\big),\qquad
(\Rs{M})^{-1}=\tfrac13\big((\hexA^{-1})_{(sr)}+2(\hexF^{-1})_{(sr)}\big),\\
\frac1{\Rs{m_3}}=\frac17\Big(c+\frac1{m_{B_1}}+\frac1{m_{B_2}}\Big),\qquad
\frac1{\Rs{m_2}}=\frac15\Big(\frac1{m_{A_1}}+2(\hexF^{-1})_{22}+2(\hexG^{-1})_{22}\Big),
\end{gathered}\label{eq:hex-voigt-reuss}
\end{equation}
with $c=(\hexA^{-1})_{33}+2(\hexF^{-1})_{33}+2(\hexG^{-1})_{33}$. The Hill estimate is
$\Hl{\tD}=\tfrac12(\Vg{\tD}+\Rs{\tD})$. In the constants $a_i$, \eqref{eq:cub-voigt}
reproduces the Voigt average of Lazar et al.~\cite[eq.~(158)]{Lazar2022}. Classical closed
forms are obtained from rotation invariants, for third-order elastic
constants~\cite[eqs.~(2.4) and~(2.6)]{Blaschke2017} and for photoelastic tensors, which
lack the major symmetry~\cite{Ranganath1973}. Kube and Norris~\cite{Kube2026} obtain the
Voigt and Reuss averages of cubic third-order constants in a basis of Walpole's type; here
the blocks give both for the gradient tensor, without invariants.

\subsection{Reuss below Voigt}\label{sec:averages-order}

The Reuss estimate needs only $\tD$ and $\avg{\tD^{-1}}$ to be invertible, but it is the
response of the grains under a uniform double stress only if $\tD\succ0$. We use it, and
the Hill estimate, only then, which excludes the Fe and W data
(section~\ref{sec:compliance-data}).

\emph{Proposition.} If $\tD\succ0$, then $\Rs{\tD}\preceq\Vg{\tD}$, that is,
$\Rs{m_3}\le\Vg{m_3}$, $\Rs{m_2}\le\Vg{m_2}$ and $\Rs{M}\preceq\Vg{M}$, with
$\Rs{\tD}=\Vg{\tD}$ if and only if $\tD$ is isotropic.

The orientation average is a positive unital linear map, so the inequality is Choi's,
$\avg{\tD}^{-1}\preceq\avg{\tD^{-1}}$~\cite[theorem~2.3.6]{Bhatia2007}; we give the short
proof because it also yields the case of equality.

\emph{Proof.} Let $\tD_Q=(Q\star)\tD(Q\star)^{\mathsf T}$. Since $\avg{\tD_Q}=\Vg{\tD}$ and
$\Vg{\tD}$ commutes with every $Q\star$, with $\mathbb V=(\Vg{\tD})^{-1}$,
\begin{equation}
\begin{aligned}
(\Rs{\tD})^{-1}-(\Vg{\tD})^{-1}
&=\avg{\tD_Q^{-1}-2\mathbb V+\mathbb V\tD_Q\mathbb V}\\
&=\avg{(\mathbb V-\tD_Q^{-1})\,\tD_Q\,(\mathbb V-\tD_Q^{-1})}\succeq0,
\end{aligned}
\label{eq:reuss-below-voigt}
\end{equation}
which is the Schur complement of the average of the positive semi-definite matrix
$\left(\begin{smallmatrix}\tD_Q&\tI\\\tI&\tD_Q^{-1}\end{smallmatrix}\right)$~\cite[theorem~1.3.3]{Bhatia2007}.
Inversion reverses the order of positive definite maps, so $\Rs{\tD}\preceq\Vg{\tD}$, and
since both are isotropic the order holds block by block. The integrand has the same trace for every
$Q$, positive unless $\tD=\Vg{\tD}$; hence equality holds only for isotropic $\tD$.
$\square$

In classical elasticity the Voigt and Reuss estimates bound the effective
stiffness~\cite{Blaschke2017,Kube2016}. Here they do not: they are the responses of the
grains under a uniform strain gradient or a uniform double stress. The heterogeneity
of $\tC$ alone induces gradient moduli that grow with the square of the grain
size~\cite{Malyi2021}, for which Voigt and Reuss bounds of the gradient effective medium
have been derived on spherical volume elements~\cite{Muhlich2012}; they are absent
from any average of $\tD$. For third-order elastic
constants, too, the bounds are lost~\cite{Blaschke2017}. We use them as algebraic
estimates of the part of the effective tensor carried by $\tD$, and the Hill mean as a
convention, empirical already in the classical case~\cite{Kube2016}.

\subsection{Cubic grains}\label{sec:averages-cubic}

By~\eqref{eq:cub-voigt} and~\eqref{eq:cub-reuss},
$\Vg{M}-\Rs{M}=\bm m_h\bm m_h^{\mathsf T}/M_{hh}$ has rank at most one: one generalised
eigenvalue of the pair $(\Rs{M},\Vg{M})$ equals one, and the two estimates agree on
$\Hrm1$ if and only if $M_{hs}=M_{hr}=0$. This continues a classical sequence. In a cubic
crystal the bulk part lies in a single irreducible block, so $\kappa_V=\kappa_R$, and
$\mu_V-\mu_R$ is proportional to the square of the anisotropy $c_{11}-c_{12}-2c_{44}$~\cite[eqs.~(27)
and~(28)]{Kube2016}. Here $\Hrm1$ lies in $\irr{T_1}$, of multiplicity three, and the
difference is of rank one. In the hexagonal class $\Hrm1$ enters $\irr{A_2}$ and $\irr{E_1}$
together with copies of $\Hrm3$ and $\Hrm2$, and the rank is two in general. A second
property follows from~\eqref{eq:cub-reuss}:
\begin{equation}
\frac{\Rs{m_2}}{m_E}=\frac{5/2}{1+\tfrac32\,m_E\,(N^{-1})_{22}}\;\to\;\frac52
\quad(m_E\to0),\label{eq:reuss-factor}
\end{equation}
where $\tfrac52=5/\dimirr{E}$: the nearly vanishing $\irr{E}$ part, two of the five
dimensions of $\Hrm2$, dominates the harmonic mean, as $\mu_R\to\tfrac52c'$ when
$c'=\tfrac12(c_{11}-c_{12})\to0$ in the classical case, by the Reuss modulus
of~\cite[eq.~(28)]{Kube2016} (their eq.~(30) prints the prefactor $\tfrac53$ for $5$).

Table~\ref{tab:averages} gives the estimates of the five data sets, and
figure~\ref{fig:cubic-averages}(a) the ratios of Reuss to Voigt. The Voigt estimates of Al
and W reproduce table~6 of Lazar et al.~\cite{Lazar2022} to $\NumTableSixMaxDiff$~eV/\AA.
For $\Hrm3$, the ratio is $\NumAlRVmthree$ (Al), $\NumCuMEAMRVmthree$ (Cu MEAM) and
$\NumCuEAMRVmthree$ (Cu EAM); for $\Hrm2$, $\NumAlRVmtwo$, $\NumCuMEAMRVmtwo^{\dagger}$ and
$\NumCuEAMRVmtwo$, the Cu MEAM value set by its small $m_E$, with
$\Rs{m_2}/m_E=\NumCuMEAMReussRatio$. The two copper tensors differ more than the Voigt and
Reuss estimates of either: the difference belongs to the potentials, not to copper.

These estimates must be weighed against the moduli induced by the heterogeneity of
$\tC$. Malyi~\cite[eqs.~(10.21) and (10.22)]{Malyi2021} gives them, to second order in the
anisotropy $c_{11}-c_{12}-2c_{44}$, for spherical grains of radius $R$ with independent
orientations, through his $L^2=2R^2/5$. Taken to form~II, they have the norm of the Voigt
estimate of $\tD$ at a grain diameter $2R$ of $\NumAlInducedCrossoverNm$~nm (Al),
$\NumCuMEAMInducedCrossoverNm$~nm (Cu MEAM) and $\NumCuEAMInducedCrossoverNm$~nm (Cu
EAM), and at $2R=100$~nm they exceed it by factors of about
$\NumAlInducedRatioHundredNm$ (Al) and $\NumCuMEAMInducedRatioHundredNm$ (Cu). This is an order of
magnitude, but it settles the scale: above a few nanometres the gradient response of a
polycrystal is set by the heterogeneity of $\tC$, and the averages of $\tD$ describe
the remainder.

For a physical reading we define block lengths $\lblk{p}$ by
\begin{equation}
\tD^p\bm w=\lblk{p}^2\,(\tC\ot\kron)^p\,\bm w,\label{eq:block-lengths}
\end{equation}
the generalised eigenvalues of each block of $\tD$ relative to that of the separable tensor
$\tC\ot\kron$; $\tD$ has the Helmholtz form $\ell^2\tC\ot\kron$ if and only if all
$\lblk{p}$ equal $\ell$. Each gradient estimate is paired with the estimate of $\tC$ of the
same kind. The pairing is a normalisation, not a joint estimate: a uniform strain and a
uniform strain gradient cannot be imposed together~\cite[remark~1]{Lazar2022}. Since $\avg{\tC\ot\kron}=\avg{\tC}\ot\kron$ and
$(\tC\ot\kron)^{-1}=\tC^{-1}\ot\kron$, a Helmholtz grain then has
$\lblk{p}=\ell$ in all three. The length of Lazar and Po~\cite[eq.~(107)]{Lazar2024} is
their dimension-weighted root mean square:
\begin{equation}
\lLP^2=\tfrac1{18}\,\big(\tC^{-1}\big)_{ijkl}D_{ijmklm}
=\tfrac1{18}\tr\big((\tC\ot\kron)^{-1}\tD\big)=\tfrac1{18}\sum_p\dimirr{p}\sum\lblk{p}^2,
\label{eq:lazar-po-length}
\end{equation}
the inner sum running over the lengths of block $p$. In the Voigt estimates
(table~\ref{tab:block-lengths}, figure~\ref{fig:cubic-averages}(b)) the block lengths
spread by factors $\NumAlVoigtLengthSpread$, $\NumCuMEAMVoigtLengthSpread$ and
$\NumCuEAMVoigtLengthSpread$, far from a single length; $\lLP$ of the Hill estimates is
$\NumAlHillLengthLP$, $\NumCuMEAMHillLengthLP$ and $\NumCuEAMHillLengthLP$~\AA. The Cu EAM
lengths of Reuss exceed those of Voigt because $\mu_R<\mu_V$ in the normalisation, without
contradicting the proposition.

\subsection{Hexagonal grains}\label{sec:averages-hexagonal}

By~\eqref{eq:hex-voigt-reuss} the Voigt estimate depends on $m_{B_1}+m_{B_2}$ only: by
\eqref{eq:hex-extra}, the extra modulus multiplies a harmonic tensor of degree six, whose
average vanishes. In the Reuss estimate only $\Rs{m_3}$ depends on it, through
$1/m_{B_1}+1/m_{B_2}=8\sigma/(4\sigma^2-\Dhex^2)$, $\sigma=\tfrac12(m_{B_1}+m_{B_2})$:
\begin{equation}
\Rs{m_3}(\Dhex)=\frac{7}{c+8\sigma/(4\sigma^2-\Dhex^2)},\label{eq:hex-reuss-delta}
\end{equation}
even in $\Dhex$, decreasing in $|\Dhex|$ and vanishing as $|\Dhex|\to2\sigma$, where
$\tP{B_1}$ or $\tP{B_2}$ loses positivity. No average depends on the sign of $\Dhex$,
that is, on the family of twofold axes chosen for $\bm a$.

No hexagonal gradient tensor of a crystal has been published, but granular
micromechanics gives one. For equal spheres with normal and tangential contact
stiffnesses $k_n$ and $k_s$, Pirmoradi et al.~\cite[eq.~(70)]{Pirmoradi2024} write the
sixth-order tensor, in form~I with the $\kappa$ of section~\ref{sec:form-II}, as an
integral over the contact directions weighted by their density, expanded in spherical
harmonics~\cite[eq.~(71)]{Pirmoradi2024}. They work out isotropic, transversely isotropic,
orthotropic, monoclinic and triclinic fabrics; the hexagonal one is our choice within their
eq.~(71). With the axial terms $a_{20}$, $a_{40}$, $a_{60}$ and the sixfold term
$a_{66}P_6^6(\cos\theta)\cos6\varphi$, where $P_6^6(x)=10395\,(1-x^2)^3$ is the
unnormalised associated Legendre function, the density is invariant under
the hexagonal group and even in the direction, so the fifth-order coupling vanishes, and
the classical tensor, which sees only the terms of degree up to four, is transversely
isotropic. In the basis, $a_{66}$ changes only $m_{B_1}$ and $m_{B_2}$:
\begin{equation}
\Dhex=\tfrac{1440}{13}\,(k_n-k_s)\,a_{66}\,\rho_cl^4,\label{eq:hex-granular}
\end{equation}
with $\rho_c$ the number of contacts per unit volume and $l$ the branch length; for
uniform contacts the tensor reproduces their eq.~(74). The extra modulus is thus the
sixfold harmonic of the contact distribution, weighted by the difference of the two
contact stiffnesses. A positive density bounds it: for $k_s=k_n/2$, $a_{20}=0.2$ and $a_{40}=a_{60}=0$, within
the range of their figure~3, $|\Dhex|\le\NumHexGranDeltaMax\,\sigma$, and $\Rs{m_3}$ changes
by at most $\NumHexGranReussChangePercent\,\%$. Sweeping $\Dhex$ further on the same base
(figure~\ref{fig:hexagonal}), $\Rs{m_3}(\sigma)/\Rs{m_3}(0)$ is $\NumHexTIReussAtSigma$, and
$21/23=\NumHexIsoReussAtSigma$ for every isotropic base, for which $c=5/m_3$ and
$\sigma=m_3$: the effect of the extra modulus is of second order.

\begin{table}[pos=tbp]
\caption{Voigt, Reuss and Hill estimates~\eqref{eq:cub-voigt}--\eqref{eq:cub-reuss} of the
cubic data sets of table~\ref{tab:block-moduli}: isotropic block moduli~\eqref{eq:iso-tensor}
(eV/\AA). Fe and W, which are not positive definite, have only the Voigt estimate. The last
digit shown is uncertain by at most five units; $\dagger$ marks a relative uncertainty
above 1\,\%, from the $m_E$ of Cu MEAM.}\label{tab:averages}
\centering\small
\input{averages}
\end{table}

\begin{table}[pos=tbp]
\caption{Block lengths~\eqref{eq:block-lengths} of the estimates (\AA), each relative to
$\tC\ot\kron$ with the classical stiffness of the same estimate, unlike
table~\ref{tab:application}, which uses the stiffness of Kr\"oner's scheme throughout;
$\lblk{1}^{\mp}$ are the two lengths of the $\Hrm1$ block, and
$\lLP$~\eqref{eq:lazar-po-length}. The row `crystal' gives $\lLP$ of the single crystal,
which for Al and Cu MEAM reproduces table~2 of Lazar and Po~\cite{Lazar2024}. Digits and
$\dagger$ as in table~\ref{tab:averages}.}\label{tab:block-lengths}
\centering\small
\input{block_lengths}
\end{table}

\begin{figure}[pos=tbp]
\centering
\includegraphics{cubic_averages}
\caption{Orientation averages of the positive definite cubic data. (a) Reuss over Voigt for
$m_3$ ($\Hrm3$), $m_2$ ($\Hrm2$) and, in the rows $\Hrm1^{\mp}$, the lower and upper
generalised eigenvalues of $(\Rs{M},\Vg{M})$; the upper one is exactly one. (b) Block
lengths of the Voigt (filled) and Reuss (open) estimates, as in
table~\ref{tab:block-lengths}, with $\lblk{1}^{\mp}$ in the rows $\Hrm1^{\mp}$, and $\lLP$ of
Voigt (dotted lines).}\label{fig:cubic-averages}
\end{figure}

\begin{figure}[pos=tbp]
\centering
\includegraphics{hexagonal_extra_modulus}
\caption{Reuss and Hill moduli $m_3$ over the Voigt modulus, against $\Dhex/\sigma$ with
$m_{B_1,B_2}=\sigma\pm\Dhex/2$, on two bases: any isotropic tensor, and the transversely
isotropic (TI) tensor of a granular packing with $k_s=k_n/2$ and
$a_{20}=0.2$~\cite{Pirmoradi2024}. Shaded: the values of $\Dhex$ that a positive contact
density gives to this packing, by~\eqref{eq:hex-granular}. Voigt does not depend on
$\Dhex$.}\label{fig:hexagonal}
\end{figure}

\section{Illustrative application: a self-consistent estimate}\label{sec:application}

The self-consistent scheme needs all four operations of the algebra on a complete cubic
tensor. It yields three results. The concentration tensor of the grain comes in blocks,
and it shows that the scheme needs six conditions where five have been imposed
(section~\ref{sec:application-blocks}). For grains of physical size the scheme has a
closed limit, in which the polycrystal sees the grain only through its Voigt estimate
(section~\ref{sec:application-limit}). And the average of the Eshelby tensor on which the
scheme rests is obtained exactly (appendix~\ref{app:eshelby}).

Solyaev~\cite{Solyaev2023} extended Kr\"oner's self-consistent scheme to strain-gradient
elasticity. A spherical grain of radius $R$, with tensors $\tC_\gr$ and $\tD_\gr$, lies in
the isotropic effective medium $(\tC,\tD)$ under a uniform strain and a uniform strain
gradient, and the averages over the grain and over its orientations must return them. He
works in form~I with the variable $\tilde\kappa_{ijk}=u_{k,ij}=\eta_{kij}+\eta_{kji}-\eta_{ijk}$,
that is, $\tilde{\bm\kappa}=(3\tP{\mathcal S}-2\tI)\sgrad$, which multiplies the copies in $\mathcal S$ by $1$ and those
in $\mathcal R$ by $-2$: a form-I block is the form-II block rescaled copy by copy, and the
algebras do not change. In this form, the gradient condition reads~\cite[eqs.~(7)
and~(10)]{Solyaev2023}
\begin{equation}
\tB=\big(\tI+\tSav\,\tD^{-1}(\tD_\gr-\tD)\big)^{-1},\qquad \avg{\tB}=\tI,
\label{eq:sc-condition}
\end{equation}
with $\tSav$ the average over the sphere of the sixth-order Eshelby tensor of the medium:
as in Solyaev's scheme, the Eshelby field, which is not uniform in the grain, is replaced
by its average. The fourth-order analogue, with $\tC$ and $\tC_\gr$, becomes Kr\"oner's condition for
large grains.

\subsection{The concentration tensor in blocks}\label{sec:application-blocks}

Since $\tD$ and $\tSav$ are isotropic, $\mathbb L=\tSav\tD^{-1}$ is an isotropic element
without major symmetry, and $\tB$ lies in $\Alg{\clscub}$ for a cubic grain, with
\begin{equation}
\tB^p=\big(I+\mathbb L^p(\tD^p_\gr-\tD^p)\big)^{-1},\label{eq:sc-blocks}
\end{equation}
where, by~\eqref{eq:cub-iso-units}, an isotropic element $(m_3,m_2,M)$ has the cubic blocks
$\mathrm{diag}(m_3,M)$ in $\irr{T_1}$, $\mathrm{diag}(m_3,m_2)$ in $\irr{T_2}$, $m_2$ in
$\irr{E}$ and $m_3$ in $\irr{A_2}$. The grain enters through one $3\times3$, one $2\times2$
and two scalar inversions. By~\eqref{eq:weighted-means}, which holds without major
symmetry, $\avg{\tB}$ has the isotropic blocks
$y_3=\tfrac17(3B^{T_1}_{hh}+3B^{T_2}_{33}+B^{A_2})$,
$y_2=\tfrac15(3B^{T_2}_{22}+2B^{E})$ and $Y=B^{T_1}_{(sr)}$, and
\begin{equation}
\avg{\tB}=\tI\iff y_3=y_2=1,\quad Y=I_2:\label{eq:sc-six}
\end{equation}
six scalar equations, as many as the dimension of $\Alg{\clsiso}$. Solyaev imposes five
linear invariants~\cite[eq.~(13)]{Solyaev2023}. The sixth is $B_2'=B_{ijjkki}$, the
transpose of his $B_2=B_{iikkjj}$, and in his variables
\begin{equation}
B_2-B_2'=3\sqrt5\,(Y_{rs}-Y_{sr}),\label{eq:sc-sixth}
\end{equation}
so the five fix $\avg{\tB}$ only if $Y$ is symmetric, which $\tSav$ does not ensure (its
$\Hrm1$ block is not symmetric; appendix~\ref{app:eshelby}). We measure the mismatch by the
residual $\scres$ of his eq.~(22), printed without squares, with every term squared: the
root sum of squares of the
relative errors of $A_1$, $A_2$ (fourth order) and $B_1,\dots,B_5$, seven terms, or eight
with $B_2'$.

\subsection{Large grains}\label{sec:application-limit}

Solyaev's medium is the Helmholtz tensor $\ell^2\tC\ot\kron$ of form~II (his eq.~(19), read
as in~\cite[eq.~(4.3)]{Ma2018} and~\cite[eq.~(3.3)]{Gao2009}; read literally with his
$\tilde\kappa=\nabla\nabla\disp$, it would put $\tC$ on the derivative indices, which is not
the Helmholtz tensor, and such readings are tried in section~S6 of the supplementary
material), taken to form~I, with the bulk modulus of the
grain, so that $A_1=3$; the unknowns are $\nu$ and $\ell$. For this medium the exact average is
$\tSav=\Phi(R/\ell)\,\hat{\tS}(\lambda,\mu)$ (appendix~\ref{app:eshelby}), with
$\Phi$ the function of~\cite[eq.~(21)]{Solyaev2023}, $\Phi\approx-\ell/(2R)$ for $R\gg\ell$. Write $\tD=\ell^2\hat\tD$.
Since the orientation average commutes with composition by isotropic elements, to first
order in $\Phi$
\begin{equation}
\avg{\tB}-\tI=-\Phi\,\hat{\tS}\hat\tD^{-1}\big(\ell^{-2}\Vg{\tD_\gr}-\hat\tD\big),
\label{eq:sc-limit}
\end{equation}
where only the Voigt estimate of the grain enters. The $A$ conditions fix $\nu$ at
Kr\"oner's value $\nu_{\mathrm K}$~\cite[eq.~(3.7)]{Blaschke2017},
\cite[eq.~(16b)]{Kube2016}, and the $B$ terms give, to leading order, $\scres=|\Phi|\,|\ell^{-2}\bm\alpha-\bm\gamma|$,
with $\bm\alpha$ and $\bm\gamma$ the invariants $B_t$ relative to their values $B_t^0$ for $\tI$, of
$\hat{\tS}\hat\tD^{-1}\Vg{\tD_\gr}$ and $\hat{\tS}$. As $|\Phi|$ grows with $\ell$, the
minimum is at
\begin{equation}
\linf^{-2}=|\bm\gamma|/|\bm\alpha|,\qquad \scres/|\Phi|\;\to\;\big|\linf^{-2}\bm\alpha-\bm\gamma\big|.
\label{eq:sc-linf}
\end{equation}

\subsection{Results}\label{sec:application-results}

Table~\ref{tab:application} gives the limit for the positive definite data. It agrees
with the full minimisation at $R=1000$~\AA, and from $2R=2$~nm on the minimising $\ell$ is
within $\NumSCNearLimitPercent$\,\% of $\linf$ (supplementary material, section~S5). Three
readings follow. First, for grains of physical size the scheme sees the grain only through
$\Vg{\tD_\gr}$: no inverse of $\tD_\gr$ enters, and the block inversions
of~\eqref{eq:sc-blocks} matter only at sub-nanometre sizes. Second, the limit diagnoses the
scheme itself. The Helmholtz medium has equal block lengths~\eqref{eq:block-lengths}, while
those of $\Vg{\tD_\gr}$, with Kr\"oner's $\tC$, spread by $\NumAlVoigtKronerSpread$,
$\NumCuMEAMVoigtKronerSpread$ and $\NumCuEAMVoigtKronerSpread$. No single $\ell$ meets the
six conditions: the normalised residual stays at
$\scres/|\Phi|=\NumAlLimResidualSeven$--$\NumCuEAMLimResidualSeven$. The residual itself
vanishes with $|\Phi|$: with seven terms it is $\NumSCResidualTwoNmMin$--$\NumSCResidualTwoNmMax$
at $2R=2$~nm and $\NumSCResidualHundredNmMin$--$\NumSCResidualHundredNmMax$ at $2R=100$~nm, so
for large grains every $\ell$ of the order of 1~\AA\ nearly meets the conditions,
and $\ell$ is only weakly determined, as the minimiser of the coefficient of a vanishing
quantity. That minimiser $\linf$, a weighted compromise, exceeds $\lLP$ of the same Voigt estimate by $\NumAlLimOverLengthPercent$,
$\NumCuMEAMLimOverLengthPercent$ and $\NumCuEAMLimOverLengthPercent$\,\%. A one-length
medium cannot represent these polycrystals, and a medium with five gradient constants is
the next step, for which appendix~\ref{app:eshelby} already gives the general average.
Third, the sixth condition $B_2'=3$, which the seven-term solutions miss by up to
$\NumSCSixthSolutionMaxPercent$\,\%, moves $\linf$ by at most $\NumCuEAMLimShiftPercent$\,\%. Like the averages, the scheme
ignores the moduli induced by the heterogeneity of $\tC$~\cite{Malyi2021}, which at the
grain sizes where it reduces to $\Vg{\tD_\gr}$ are the dominant ones
(section~\ref{sec:averages-cubic}): it illustrates the algebra, not a homogenisation
theory.

With the grain constants of Solyaev's iron example, and grain and medium both written as the
form~II Helmholtz model, the scheme returns $\ell=\NumFeAllEllMin$--$\NumFeAllEllMax\,\ell_0$
and Kr\"oner's $\nu$ over the grain sizes of his figure~2, with any of the three averaged
tensors (the exact one, that of Ma et al.\ and his printed one); we could not reproduce the
published increase of $\ell$ (supplementary material, section~S6).

\begin{table}[pos=tbp]
\caption{Self-consistent estimate for large grains~\eqref{eq:sc-linf}: Poisson's ratio
$\nu_{\mathrm K}$ of Kr\"oner's classical estimate, with the bulk modulus of the grain; length
$\linf$ and normalised residual $\scres/|\Phi|$ with seven terms and, in brackets, eight;
spread $\ell_{\max}/\ell_{\min}$ of the block lengths~\eqref{eq:block-lengths} of the Voigt
estimate and its $\lLP$~\eqref{eq:lazar-po-length}, both relative to Kr\"oner's $\tC$.
Digits as in table~\ref{tab:averages}.}\label{tab:application}
\centering\small
\input{application}
\end{table}

\section{Discussion and conclusions}\label{sec:conclusions}

The bases of sections~\ref{sec:cubic} and~\ref{sec:hexagonal} turn every operation on
cubic and hexagonal gradient tensors into operations on blocks of size at most three and
four. The compliance, the three estimates and the concentration tensor then follow in closed
form. The blocks also read the data: they place the indefiniteness of the Fe and W tensors,
the negative root of the quadratic factor of Lazar et al., in the $\irr{T_2}$ block, the near singularity
of the Cu MEAM tensor in $m_E$, and the anisotropy of the gradient response in the spread
of the block lengths.

The hexagonal extra modulus has a dynamic reading as well. For a plane wave of direction
$\bm\xi$, the gradient term of the acoustic tensor is proportional to
$Q_{il}=D_{ijmlkn}\,\xi_j\xi_m\xi_k\xi_n$~\cite[eq.~(11)]{Rosi2016}. All the hexagonal
units except $\tP{B_1}$ and $\tP{B_2}$ give a transversely isotropic contribution, while
by~\eqref{eq:hex-extra} the extra modulus adds
$\tfrac18\Dhex\,\mathrm{Re}\big((\zhex\cdot\bm\xi)^4\,\zhex\ot\zhex\big)$, whose longitudinal
component is $\tfrac18\Dhex\sin^6\theta\cos6\varphi$ for a direction at polar angle
$\theta$ from $\bm n$ and azimuth $\varphi$ from $\bm a$. Since the classical tensor is transversely isotropic, and so is a hexagonal micro-inertia
tensor, of order four, $\Dhex$ is, without fifth-order coupling, the only modulus through which dispersive
waves feel the sixfold axis; in a granular packing it is the sixfold harmonic of the
contact distribution~\eqref{eq:hex-granular}. This is the three-dimensional
counterpart of the anisotropic dispersion of hexagonal lattices in two
dimensions~\cite{Rosi2016}.

The scope has limits. Lattices with a basis, such as hcp and diamond, are excluded: where
their gradient tensor has the full class, at an atom site, the fifth-order coupling does not
vanish (section~\ref{sec:algebras}). The Reuss and Hill estimates are used only for
positive definite tensors, which excludes the Fe and W data, and they are estimates, not
bounds; none contains the moduli induced by the heterogeneity of $\tC$~\cite{Malyi2021},
which by Malyi's estimate dominate from grain diameters of a few nanometres (Al) or below
one nanometre (Cu) on.
The hexagonal tensors are those of a granular model, not of a crystal.

The construction extends to the other classes; groups with irreducible representations
of complex type, such as the cyclic ones, need complex blocks. The fifth-order coupling,
which governs the dispersion at first order in the wavenumber in non-centrosymmetric and
chiral cubic architectured materials~\cite{Rosi2024}, enlarges the algebra to the whole space of strains and strain
gradients. A self-consistent medium with five gradient constants needs no new Eshelby tensor:
appendix~\ref{app:eshelby} covers the general isotropic medium. With the blocks, the
log-Euclidean anisotropy index of Kube and de
Jong~\cite[eq.~(20)]{Kube2016} has a direct gradient analogue. Po et
al.~\cite{Po2018} computed the complete gradient tensor of hcp magnesium but published only
a projection of it; with its fifth-order coupling, such data would test the hexagonal class
on a crystal.


\appendix
% Elsevier numbers the equations of appendix A as (A.1), (A.2), ...
\counterwithin{equation}{section}
\renewcommand{\theequation}{\thesection.\arabic{equation}}
\section{Conversion to the copies of Lazar (2016)}\label{app:lazar}

Lazar~\cite{Lazar2016} splits $\sgrad$ first into the dilatation gradient $\eta_{ppk}$ and
the shear gradient, and obtains two vector pieces: the dilatation-gradient piece
${}^{(4)}\eta_{ijk}=\tfrac13\delta_{ij}\eta_{ppk}$ and the shear-divergence piece
${}^{(3)}\eta_{ijk}$, built from the trace of the shear gradient~\cite[eqs.~(53)--(54)]{Lazar2016}.
Their unit embeddings, $(\emb{(4)}v)_{ijk}=\delta_{ij}v_k/\sqrt3$ and $\emb{(3)}v$, a
positive multiple of $\delta_{ik}v_j+\delta_{jk}v_i-\tfrac23\delta_{ij}v_k$, span the same
plane as $\emb{s}$ and $\emb{r}$:
\begin{equation}
(\emb{(4)},\emb{(3)})=(\emb{s},\emb{r})\,R_\theta,\qquad
R_\theta=\begin{pmatrix}\cos\theta&\sin\theta\\ \sin\theta&-\cos\theta\end{pmatrix},\qquad
\cos\theta=\tfrac{\sqrt5}{3},\quad \sin\theta=\tfrac23.\label{eq:lazar-pair}
\end{equation}
$R_\theta$ is a reflection: the rotation by $\theta=\arccos(\sqrt5/3)\approx41.81^\circ$
followed by the reversal of piece~(3). The irreducible gradient parameters $b_1,\dots,b_5$
of~\cite[eq.~(60)]{Lazar2016} are $b_1=m_3$, $b_2=m_2$ and
\begin{equation}
M=R_\theta\begin{pmatrix}b_4&\tfrac{\sqrt5}{3}b_5\\ \tfrac{\sqrt5}{3}b_5&b_3\end{pmatrix}R_\theta,
\qquad\text{that is,}\quad
\begin{aligned}
M_{ss}&=\tfrac19(5b_4+4b_3)+\tfrac{20}{27}b_5,\\
M_{sr}&=\tfrac{2\sqrt5}{9}(b_4-b_3)-\tfrac{\sqrt5}{27}b_5,\\
M_{rr}&=\tfrac19(4b_4+5b_3)-\tfrac{20}{27}b_5,
\end{aligned}\label{eq:lazar-block}
\end{equation}
where $\sqrt5/3$ is the value of his cross term
${}^{(3)}\eta_{ijk}\,{}^{(4)}\eta_{kji}$ for unit pieces. In the cubic class, $s$ and $r$
are copies of $\irr{T_1}$: \eqref{eq:lazar-block} converts the $(s,r)$ part of $M$, and
the couplings $(M_{hs},M_{hr})$ transform by $R_\theta$. In the hexagonal class the same
holds for the $(s,r)$ parts of the $\irr{A_2}$ and $\irr{E_1}$ blocks. The change is
orthogonal within each block, so inverses, eigenvalues and positive definiteness are
unaffected (section~\ref{sec:algebras}).

Neither pair is singled out by $\Orth$: Lazar's is adapted to the traces of $\sgrad$, the
$\GL$ pair to the splitting $\Vsg=\mathcal S\oplus\mathcal R$. The Helmholtz-type separable
tensor $\ell^2\tC\ot\kron$, with isotropic $\tC$ of Lamé constants $\lambda$ and $\mu$, is
diagonal in Lazar's pair, $b_4=(3\lambda+2\mu)\ell^2$, $b_3=2\mu\ell^2$, $b_5=0$, and
coupled in the $\GL$ pair, $M_{sr}=\tfrac{2\sqrt5}{3}\lambda\ell^2$; a couple-stress
material is uncoupled in the $\GL$ pair (section~\ref{sec:isotropic}) and coupled in
Lazar's.

\section{The averaged Eshelby tensor of the sphere}\label{app:eshelby}

We use the kinematics and constants of Ma et al.~\cite{Ma2018}: form~I with
$\tilde\kappa_{ijk}=u_{k,ij}$, the isotropic gradient tensor $\tD$ of their eq.~(2.6), with
constants $a_1,\dots,a_5$, and the lengths $l_1$, $l_2$ of their eq.~(2.8). Let a uniform
eigen-gradient $\tilde{\bm\kappa}^*$ act in a sphere $\Omega$ of radius $R$, with indicator
$\chi$. The medium has the acoustic tensor
\begin{equation}
K(\bm\xi)=\xi^2\big[(\lambda+2\mu)(1+l_1^2\xi^2)\,\bm n\ot\bm n
+\mu(1+l_2^2\xi^2)(\kron-\bm n\ot\bm n)\big],\qquad \bm n=\bm\xi/\xi,
\label{eq:esh-acoustic}
\end{equation}
and, solving in Fourier space, the average of $\tilde{\bm\kappa}$ over $\Omega$ is
$\tSav\tilde{\bm\kappa}^*$, with
\begin{equation}
\bar S_{rts,lmn}=\frac{1}{|\Omega|\,(2\pi)^3}\int|\hat\chi(\bm\xi)|^2\,
\xi_r\xi_t\,(K^{-1})_{sk}\,\xi_i\xi_j\,D_{ijk,lmn}\,\dd^3\xi,
\label{eq:esh-fourier}
\end{equation}
the indices $(r,t)$ and $(i,j)$ being those of the derivatives. Since
$\tilde\kappa$ is symmetric in its derivative indices, $\tSav$ is compared with other
expressions after symmetrisation in $(r,t)$ and in $(l,m)$. The integral separates into a
radial and an angular part:
\begin{equation}
\tSav=\mathbb P\,\tD,\qquad
P_{rts,ijk}=\frac{F(l_1)}{\lambda+2\mu}\,\avg{n_rn_tn_sn_in_jn_k}
+\frac{F(l_2)}{\mu}\big(\avg{n_rn_tn_in_j}\delta_{sk}-\avg{n_rn_tn_sn_in_jn_k}\big),
\label{eq:esh-polarization}
\end{equation}
with $\avg{\cdot}$ the average over the unit sphere and
\begin{equation}
F(l)=\frac{3}{2R^3l}\big[R^2(1+\beta)+2lR\beta-l^2(1-\beta)\big],\qquad \beta=e^{-2R/l}.
\label{eq:esh-radial}
\end{equation}
$\mathbb P$ is symmetric, a gradient analogue of Hill's polarisation tensor, and depends on the
medium only through $\lambda$, $\mu$, $l_1$ and $l_2$; $\tSav$ is not symmetric. As
$R/l_1,R/l_2\to\infty$, $l^2F\to0$ and $\tSav\to0$; as $R/l_1,R/l_2\to0$, $l^2F\to1$, the
integrand tends to a projection of rank three for every $\bm\xi$, and $\tr\tSav\to3$. For the
Helmholtz medium $\tD=\ell^2\tC\ot\kron$, that is, $a_1=\tfrac12\lambda\ell^2$, $a_2=a_3=0$,
$a_4=a_5=\tfrac12\mu\ell^2$~\cite[eq.~(4.3)]{Ma2018} and $l_1=l_2=\ell$, one has
$\ell^2F(\ell)=-3\Phi(R/\ell)$ with $\Phi$ of~\cite[eq.~(21)]{Solyaev2023}, and the blocks
of $\tSav$ in the copies of section~\ref{sec:isotropic} are
\begin{equation}
\begin{gathered}
m_3=-\frac{4(4\lambda+11\mu)}{35(\lambda+2\mu)}\,\Phi,\qquad m_2=-\frac{\Phi}{5},\qquad
M=-\Phi\begin{pmatrix}\dfrac{19\lambda+26\mu}{15(\lambda+2\mu)}&
-\dfrac{2\sqrt5\,(\lambda-\mu)}{15(\lambda+2\mu)}\\[1.5ex]
\dfrac{4\sqrt5}{15}&\dfrac13\end{pmatrix},
\end{gathered}\label{eq:esh-helmholtz}
\end{equation}
which define $\hat{\tS}=\tSav/\Phi$ of section~\ref{sec:application}. The fourth-order
average, $s_1$ and $s_2$ of Ma et al.~\cite[eq.~(3.19)]{Ma2018}, follows in the same way and
agrees with theirs.

Their sixth-order average~\cite[eq.~(3.18)]{Ma2018}, written as five isotropic tensors with
coefficients $\hat n_1,\dots,\hat n_5$, is not $\tSav$. An isotropic map of $\Vsg$ without major
symmetry has six independent coefficients ($\Alg{\clsiso}$ has dimension six), and the term
$\hat n_2$ lumps two classes that are transposes of each other,
$\delta_{lm}(\delta_{nr}\delta_{ts}+\delta_{nt}\delta_{rs})$ and
$\delta_{rt}(\delta_{sl}\delta_{mn}+\delta_{sm}\delta_{nl})$, whose coefficients in $\tSav$
differ. Of the five coefficients only $\hat n_1$ is that of $\tSav$; $\hat n_4$ and $\hat n_5$
are twice and four times those of $\tSav$. For a general isotropic medium, with $c=\lambda+2\mu$, $F_1=F(l_1)$ and
$F_2=F(l_2)$, the exact average is
$\tSav=\bar n_1\mathbb K_1+\bar n_{2a}\mathbb K_{2a}+\bar n_{2b}\mathbb K_{2b}+\bar n_3\mathbb K_3
+\bar n_4\mathbb K_4+\bar n_5\mathbb K_5$, in the classes of their eq.~(3.18) with the term of $\hat n_2$ split in two (defined below), with
\begin{equation}
\begin{aligned}
\bar n_1&=\frac{(7a_2+14a_3+2a_4+2a_5)F_1}{105c}+\frac{2(28a_3+6a_4-a_5)F_2}{105\mu},\\
\bar n_{2a}&=\frac{(14a_1+7a_2+4a_4+4a_5)F_1}{210c}+\frac{(14a_2-2a_4+5a_5)F_2}{105\mu},\\
\bar n_{2b}&=\frac{(7a_2+14a_3+2a_4+2a_5)F_1}{105c}-\frac{2(7a_3+a_4+a_5)F_2}{105\mu},\\
\bar n_3&=\frac{2(a_4+a_5)F_1}{105c}-\frac{(2a_4-5a_5)F_2}{105\mu},\qquad
\bar n_4=\frac{2(a_4+a_5)F_1}{105c}+\frac{2(6a_4-a_5)F_2}{105\mu},\\
\bar n_5&=\frac{(14a_1+7a_2+4a_4+4a_5)F_1}{210c}-\frac{(7a_2+4a_4+4a_5)F_2}{210\mu},
\end{aligned}\label{eq:esh-six-general}
\end{equation}
the exact counterparts of $\hat n_1,\dots,\hat n_5$ of~\cite[eq.~(3.18)]{Ma2018}. For the
Helmholtz medium, $l_1=l_2=\ell$ and $\ell^2F=-3\Phi$, so that
\begin{equation}
\begin{aligned}
\tSav={}&-\frac{(5\lambda+12\mu)\Phi}{35c}\big(\mathbb K_1+\mathbb K_4\big)
-\frac{(\lambda+\mu)\Phi}{7c}\,\mathbb K_{2a}+\frac{2(\lambda+\mu)\Phi}{35c}\,\mathbb K_{2b}\\
&-\frac{(3\lambda+10\mu)\Phi}{70c}\,\mathbb K_3-\frac{(3\lambda-4\mu)\Phi}{70c}\,\mathbb K_5,
\end{aligned}
\end{equation}
where $\mathbb K_1$, $\mathbb K_3$, $\mathbb K_4$ and $\mathbb K_5$ are the tensors multiplying
$\hat n_1$, $\hat n_3$, $\hat n_4$ and $\hat n_5$ in~\cite[eq.~(3.18)]{Ma2018}, and
$\mathbb K_{2a}=\delta_{rt}(\delta_{sl}\delta_{mn}+\delta_{sm}\delta_{nl})$ and
$\mathbb K_{2b}=\delta_{lm}(\delta_{nr}\delta_{ts}+\delta_{nt}\delta_{rs})$ are the two halves
of the tensor multiplying $\hat n_2$. Their eq.~(3.18) needs $\mathbb K_{2a}$ and $\mathbb K_{2b}$
to share a coefficient; for the Helmholtz medium the coefficient of $\mathbb K_{2b}$ is $-2/5$
times that of $\mathbb K_{2a}$.

The tensor of their eq.~(3.18) times $\tD^{-1}$ is not symmetric, and
for the Helmholtz medium its trace tends to $11(31\lambda+56\mu)/(35(\lambda+2\mu))$ instead of
three as $R/\ell\to0$. Their own derivation is right up to the last step: the tensor of their
eq.~(3.12), averaged over the sphere with their eq.~(B.12), gives $\tSav$ exactly, and the error
enters when that average is collected into the five terms of their eq.~(3.18). A real-space
route from their Green function confirms $\tSav$ (supplementary material, section~S7). The
tensor of Solyaev~\cite[eq.~(21)]{Solyaev2023} takes its coefficients from eq.~(3.18) with
eq.~(4.3); it agrees with $\tSav$ on two of the six classes, and its trace tends to
$2(74\lambda+125\mu)/(35(\lambda+2\mu))$. His fourth-order average~\cite[eq.~(20)]{Solyaev2023}
has half the deviatoric term of theirs. Section~\ref{sec:application} uses $\tSav$.


% Back matter for the European Journal of Mechanics A/Solids (Elsevier): CRediT, competing
% interests, funding, data availability and the declaration of generative AI, as unnumbered
% sections before the references. The declaration of competing interest is also entered in
% the declarations tool of the submission system.

\section*{CRediT authorship contribution statement}

\textbf{Daniel Gonz\'alez-Monsalve:} Conceptualization, Methodology, Formal analysis,
Software, Validation, Investigation, Data curation, Visualization, Writing -- original draft,
Writing -- review \& editing.

\section*{Declaration of competing interest}

The author declares that he has no known competing financial interests or personal
relationships that could have appeared to influence the work reported in this paper.

\section*{Funding}

This research did not receive any specific grant from funding agencies in the public,
commercial, or not-for-profit sectors.

\section*{Data availability}

The code and data that produce every table, figure and computed number of this paper and
of its supplementary material are available at
\url{https://github.com/daniel-gonzalezm/gradinv} and archived at Zenodo. The repository
contains the Python package that implements the bases, inverses, averages and
concentration tensors; the input data with the provenance of every number
(supplementary material, section~S3); the scripts that write the tables, figures and
numbers; and the test suite, whose key tests are listed in section~S2 of the
supplementary material. The atomistic constants used as input are those published by
Lazar et al.~\cite{Lazar2022} and Po et al.~\cite{Po2019}; no new data were generated.
The code is released under the MIT licence; the data, tables and figures under the
Creative Commons Attribution 4.0 licence.\pending{Zenodo DOI and its reference entry.}

\section*{Declaration of generative AI and AI-assisted technologies in the manuscript preparation process}

During the preparation of this work the author used Claude (Anthropic) to check
derivations, write and test code, search the literature, and draft and edit the text. After
using this tool, the author reviewed and edited the content as needed and takes full
responsibility for the content of the published article.


\bibliographystyle{unsrtnat}
\bibliography{references}

\end{document}
